\documentclass[11pt,a4paper]{article}

\usepackage[T1]{fontenc}
\usepackage[utf8]{inputenc}
\usepackage[italian]{babel}
\usepackage[a4paper,margin=2.2cm]{geometry}
\usepackage{amsmath,amssymb}
\usepackage{array}
\usepackage{booktabs}
\usepackage{longtable}
\usepackage{tabularx}
\usepackage{graphicx}
\usepackage{xcolor}
\usepackage{hyperref}
\usepackage{enumitem}
\usepackage{listings}
\usepackage{tikz}
\usepackage{pgfplots}
\usetikzlibrary{arrows.meta,calc,fit,positioning,shapes.geometric,shapes.misc}
\pgfplotsset{compat=1.18}
\usepackage[strings]{underscore}
\hypersetup{
  colorlinks=true,
  linkcolor=blue!50!black,
  urlcolor=blue!50!black
}
\urlstyle{same}

\definecolor{codebg}{RGB}{247,248,250}
\definecolor{factblue}{RGB}{28,82,181}
\definecolor{warnred}{RGB}{170,30,45}
\definecolor{okgreen}{RGB}{32,120,74}
\definecolor{softgray}{RGB}{90,90,90}

\lstdefinestyle{shellstyle}{
  basicstyle=\ttfamily\small,
  backgroundcolor=\color{codebg},
  frame=single,
  rulecolor=\color{black!15},
  breaklines=true,
  columns=fullflexible
}

\lstdefinestyle{cppstyle}{
  language=C++,
  basicstyle=\ttfamily\small,
  backgroundcolor=\color{codebg},
  frame=single,
  rulecolor=\color{black!15},
  breaklines=true,
  columns=fullflexible,
  keywordstyle=\color{factblue}\bfseries,
  commentstyle=\color{softgray},
  stringstyle=\color{okgreen!60!black},
  showstringspaces=false
}

\newcommand{\pkg}{\texttt{auto\_aim\_targeting}}
\newcommand{\code}[1]{\texttt{#1}}
\newcommand{\topic}[1]{\texttt{#1}}

\title{\pkg: guida tecnica completa\\
\large Dalle fondamenta al livello avanzato}
\author{Documento derivato dall'analisi diretta del codice del repository}
\date{Snapshot del workspace: 18 aprile 2026}

\begin{document}
\maketitle
\tableofcontents
\clearpage

\section{Breve conclusione}

Il package \pkg\ e' un runtime ROS~2 che unifica in un solo nodo quattro lavori che, concettualmente, appartengono a tutta la catena di auto-aim:
\begin{enumerate}[leftmargin=1.4em]
  \item ricevere detection 2D di armature da un detector;
  \item ricostruire misure 3D consistenti nel frame di tracking;
  \item mantenere uno o piu' target stimati tramite tracking probabilistico;
  \item decidere dove puntare, quando sparare e quali messaggi pubblicare.
\end{enumerate}

La cosa piu' importante da capire e' questa: \emph{il package non insegue semplicemente una bounding box}. Insegue un modello geometrico e dinamico del robot avversario. Per questo tiene traccia del centro del bersaglio, della sua velocita', della sua velocita' angolare, del raggio delle armature rispetto al centro e perfino del disallineamento verticale tra coppie di facce. Da qui derivano tre capacita' avanzate:
\begin{enumerate}[leftmargin=1.4em]
  \item sopravvivere a occlusioni brevi;
  \item gestire il cambio di faccia durante la rotazione del robot nemico;
  \item pianificare tiri diretti o indiretti in base a spin, latenza e balistica.
\end{enumerate}

La lettura corretta del package e' quindi: \emph{measurement -> tracking -> planning -> fire gating -> output}, non ``detector -> comando''. Se si capisce questo, si capiscono anche i parametri, i topic e gli errori di integrazione piu' comuni.

\subsection{Come usare questa guida se parti da zero}

Questa guida e' lunga perche' il package fa molte cose insieme. Un principiante non dovrebbe provare a memorizzare tutti i parametri al primo passaggio. L'ordine consigliato e':
\begin{enumerate}[leftmargin=1.4em]
  \item capire il vocabolario minimo: detection, PnP, TF, tracker, planner, fire gate;
  \item seguire il percorso di un singolo frame, dalla camera fino al comando gimbal;
  \item leggere prima i tipi dati principali, poi i singoli algoritmi;
  \item usare i test come esempi piccoli e controllati;
  \item solo alla fine ragionare sul tuning dei parametri.
\end{enumerate}

Una regola pratica utile: quando vedi un nome di classe, chiediti sempre \emph{che cosa entra, che cosa esce, e quale stato interno viene modificato}. Per esempio:
\begin{itemize}[leftmargin=1.4em]
  \item \code{DetectionConverter}: entra una detection 2D, esce una lista di armature 3D;
  \item \code{Tracker}: entra una misura 3D nel tempo, esce uno stato stimato e predicibile;
  \item \code{ShotPlanner}: entra uno snapshot del tracker, esce un possibile piano di tiro;
  \item \code{EngagementPlanner}: entrano piu' piani possibili, esce il migliore;
  \item \code{FireGate}: entra il piano migliore, esce il permesso o il blocco di fuoco;
  \item \code{GimbalCommandController}: entra il piano autorizzato, escono i messaggi di comando.
\end{itemize}

Se qualcosa non funziona, non saltare subito ai parametri del planner. Quasi sempre gli errori piu' gravi nascono prima: calibrazione camera, TF, timestamp, pose source, classi del detector o geometria PnP.

\section{Spiegazione dettagliata: dalle fondamenta al livello avanzato}

\subsection{Che problema risolve davvero questo package}

Un sistema di auto-aim deve rispondere a una domanda molto concreta:
\begin{quote}
  dato cio' che la camera vede ora, quale comando devo inviare al gimbal adesso perche' il proiettile colpisca un'armatura reale tra qualche decina o centinaio di millisecondi?
\end{quote}

Questa domanda e' difficile per almeno cinque motivi:
\begin{enumerate}[leftmargin=1.4em]
  \item la camera osserva immagini 2D, ma il tiro richiede geometria 3D;
  \item i detector producono misure rumorose, discontinue e talvolta false;
  \item il bersaglio puo' traslare e ruotare;
  \item il proiettile impiega tempo a viaggiare e cade per gravita';
  \item il sistema ha latenza: immagine, trasporto, elaborazione, attuazione.
\end{enumerate}

\pkg\ esiste per decomporre questo problema in sotto-problemi ingegneristicamente gestibili.

\subsection{Responsabilita' del package a livello di architettura}

Dal punto di vista concettuale, \pkg\ fa il lavoro che in un'architettura piu' vecchia era separato in un tracker e in un solver di traiettoria. Il README del package lo dichiara esplicitamente: il package consuma detector output e stato camera, produce target tracciati e comandi gimbal, e mantiene stabili le interfacce ROS pubbliche mentre consolida la pipeline interna.

L'entrypoint canonico del codice e':
\begin{itemize}[leftmargin=1.4em]
  \item componente ROS~2: \code{rm\_auto\_aim::AutoAimNode};
  \item file sorgente: \path{src/auto_aim_targeting/src/auto_aim_node.cpp};
  \item launch dichiarato come ``owner-path'' nel README: \path{src/launch_pkg/launch/debug_targeting.launch.py}.
\end{itemize}

Nel costruttore del nodo si vedono tre fasi:
\begin{enumerate}[leftmargin=1.4em]
  \item caricamento configurazione;
  \item creazione interfacce ROS;
  \item costruzione pipeline interna.
\end{enumerate}

Questo e' un segnale architetturale utile: la pipeline e' composta da oggetti specializzati e il nodo principale fa soprattutto orchestrazione.

\subsection{Concetti fondamentali da capire prima dei dettagli}

\subsubsection{Detection 2D}

Il detector fornisce un messaggio \code{vision\_msgs::msg::Detection2DArray} sul topic \topic{/detector/armors}. Ogni detection contiene, tra le altre cose, una bounding box 2D e una lista di classificazioni candidate.

Il package non usa direttamente la bounding box come bersaglio finale. La bounding box e' solo un indizio grezzo da cui costruire una misura geometrica.

\subsubsection{PnP: da immagine a posa 3D}

PnP significa \emph{Perspective-n-Point}. In pratica:
\begin{equation}
  \text{punti 3D noti sul modello} + \text{punti 2D osservati nell'immagine} + \text{intrinseci camera}
  \Longrightarrow
  (R, t)
\end{equation}

dove \(R\) e \(t\) sono rotazione e traslazione dell'oggetto rispetto alla camera.

Nel package:
\begin{itemize}[leftmargin=1.4em]
  \item il modello 3D e' una piastra di armatura RoboMaster, piccola o grande;
  \item i punti 2D sono stimati sinteticamente dai lati della bounding box;
  \item gli intrinseci vengono da \topic{/camera_info}.
\end{itemize}

\subsubsection{Frame di riferimento}

Un punto 3D senza frame e' ambiguo. Il package usa almeno due frame concettuali:
\begin{enumerate}[leftmargin=1.4em]
  \item il frame ottico della camera, tipicamente \code{camera\_color\_optical\_frame};
  \item il frame di tracking, configurabile tramite \code{target\_frame}, con default \code{odom}.
\end{enumerate}

La misura nasce nel frame camera, poi viene trasformata nel frame di tracking tramite TF2:
\begin{equation}
  \mathbf{p}_{\text{tracking}} = T_{\text{tracking} \leftarrow \text{camera}} \, \mathbf{p}_{\text{camera}}
\end{equation}

\subsubsection{Gimbal pose vs target pose}

Il package tiene separati due concetti:
\begin{itemize}[leftmargin=1.4em]
  \item la posa del bersaglio nel mondo;
  \item la posa attuale del gimbal/camera rispetto allo stesso frame.
\end{itemize}

Questa separazione serve per calcolare \emph{angoli relativi} di tiro e verificare se il gimbal puo' fisicamente raggiungere in tempo la soluzione proposta.

\subsubsection{Tracking probabilistico}

Una singola misura non basta. Il package usa un \emph{Extended Kalman Filter} (EKF) per stimare lo stato del target nel tempo, filtrare rumore, prevedere il futuro vicino e mantenere coerenza quando la misura salta o sparisce per pochi frame.

\subsubsection{Planning di tiro}

Dopo il tracking, il problema cambia: non bisogna piu' chiedersi ``dove si trova il target adesso?'', ma ``quale faccia devo ingaggiare e con quale anticipo temporale?''.

Questa e' la responsabilita' del planner:
\begin{itemize}[leftmargin=1.4em]
  \item scegliere il target migliore tra piu' tracker;
  \item scegliere la faccia migliore del target;
  \item decidere se usare tiro diretto visibile, tiro diretto predetto o tiro indiretto;
  \item passare la soluzione a un fire gate che decide se sparare davvero.
\end{itemize}

\subsection{Ingressi, uscite, topic e servizio}

\begin{longtable}{p{0.25\textwidth} p{0.28\textwidth} p{0.39\textwidth}}
\toprule
\textbf{Interfaccia} & \textbf{Tipo} & \textbf{Ruolo operativo} \\
\midrule
\endfirsthead
\toprule
\textbf{Interfaccia} & \textbf{Tipo} & \textbf{Ruolo operativo} \\
\midrule
\endhead
\topic{/detector/armors} & \code{vision\_msgs::msg::Detection2DArray} & Ingresso principale: detection 2D delle armature. Triggera il callback per-frame. \\
\topic{/camera_info} & \code{sensor\_msgs::msg::CameraInfo} & Fornisce intrinseci e distorsione. Serve a inizializzare il solver PnP. \\
\topic{/micro_pose} & \code{geometry\_msgs::msg::PoseStamped} & Sorgente opzionale di pitch/yaw del gimbal quando \code{pose\_source=micro\_pose}. \\
\topic{/imu/data} & \code{sensor\_msgs::msg::Imu} & Sorgente opzionale di orientazione quando \code{pose\_source=camera\_imu}. \\
\topic{/tracker/target} & \code{auto\_aim\_interfaces::msg::Target} & Miglior target selezionato o target vuoto se nessun best target. \\
\topic{/tracker/targets} & \code{auto\_aim\_interfaces::msg::Targets} & Collezione di tutti i target non-LOST e indice del best target. \\
\topic{/tracker/info} & \code{auto\_aim\_interfaces::msg::TrackerInfo} & Diagnostica di tracking: errori, yaw diff, face angle, damping effettivo. \\
\topic{/detections_output/optimal_target} & \code{vision\_msgs::msg::Detection2D} & Bounding box originale associata al best tracker, se disponibile. \\
\topic{/tracker/cmd_gimbal} & \code{auto\_aim\_interfaces::msg::GimbalCmd} & Comando finale di pitch/yaw e fire flag. \\
\topic{/cmd_vel} & \code{geometry\_msgs::msg::Twist} & Canale di compatibilita': nel codice corrente porta pitch/yaw assoluti in radianti su \code{angular.y/z}; non e' una twist fisica standard. \\
\topic{/trajectory/marker} & \code{visualization\_msgs::msg::Marker} & Marker del punto di impatto previsto. \\
\topic{/tracker/marker} & \code{visualization\_msgs::msg::MarkerArray} & Marker RViz di centro, velocita' e armature. \\
\topic{/tracker/reset} & \code{std\_srvs::srv::Trigger} & Reset esplicito di tutti i tracker e dello stato di selezione corrente. \\
\bottomrule
\end{longtable}

\subsection{Flusso completo per frame}

La sequenza reale per ogni frame di detection e' molto vicina a questa:

\begin{figure}[ht]
\centering
\begin{tikzpicture}[
  node distance=6mm and 5mm,
  box/.style={draw, rounded corners=2pt, align=center, minimum width=2.9cm, minimum height=0.95cm, fill=blue!6},
  tinybox/.style={draw, rounded corners=2pt, align=center, minimum width=2.5cm, minimum height=0.9cm, fill=green!7},
  warnbox/.style={draw, rounded corners=2pt, align=center, minimum width=2.6cm, minimum height=0.9cm, fill=orange!10},
  arr/.style={-{Latex[length=2.5mm]}, thick}
]
  \node[box] (det) {Detection2DArray\\\topic{/detector/armors}};
  \node[box, right=of det] (pose) {GimbalPoseAdapter\\freshness + TF};
  \node[box, right=of pose] (meas) {DetectionConverter\\2D -> 3D + TF};
  \node[box, right=of meas] (trk) {Trackers\\predict / assign / update / spawn / prune};
  \node[box, below=of trk] (snap) {TrackSnapshot[]};
  \node[box, left=of snap] (lat) {LatencyCompensator};
  \node[box, left=of lat] (ctx) {PlanningContext};
  \node[box, below=of lat] (plan) {ShotPlanner +\\EngagementPlanner};
  \node[warnbox, right=of plan] (fire) {FireGate};
  \node[tinybox, right=of fire] (out) {OutputPublisher\\Target + Cmd + Marker};

  \draw[arr] (det) -- (pose);
  \draw[arr] (pose) -- (meas);
  \draw[arr] (meas) -- (trk);
  \draw[arr] (trk) -- (snap);
  \draw[arr] (snap) -- (lat);
  \draw[arr] (lat) -- (ctx);
  \draw[arr] (ctx) -- (plan);
  \draw[arr] (plan) -- (fire);
  \draw[arr] (fire) -- (out);
  \draw[arr] (snap.south) |- (plan.east);
\end{tikzpicture}
\caption{Pipeline concettuale di \pkg\ per un frame di detection.}
\end{figure}

Nel codice, la pipeline entra da \code{AutoAimNode::armorsCallback()} e poi chiama \code{processFrame()}. Dentro \code{processFrame()} si vedono, nell'ordine, questi passi:
\begin{enumerate}[leftmargin=1.4em]
  \item protezione dello stato tracker con mutex;
  \item costruzione di input vuoti sicuri se qualche puntatore e' nullo;
  \item aggiornamento del \code{dt};
  \item prediction di tutti i tracker attivi;
  \item assegnazione detection-tracker;
  \item update dei tracker con detection compatibili;
  \item spawn di nuovi tracker da detection non assegnate;
  \item pruning dei tracker persi;
  \item raccolta di snapshot immutabili;
  \item aggiornamento latenza e costruzione del planning context;
  \item selezione del best plan;
  \item valutazione del fire gate;
  \item costruzione e pubblicazione dell'output finale.
\end{enumerate}

Esiste anche un watchdog sul detector: se i frame smettono di arrivare oltre \code{detector\_stall\_timeout}, il nodo pubblica un hold output. Questo evita di lasciare in uscita un comando vecchio quando il detector si ferma.

\subsection{Sottosistema di misura: da bounding box a armatura 3D}

\subsubsection{Passo 1: camera intrinsics}

Il solver PnP viene creato solo quando arriva \topic{/camera_info}. Questo avviene in \code{cameraInfoCallback()}. Un dettaglio importante: dopo aver ricevuto il primo \code{CameraInfo}, il nodo distrugge la subscription a \topic{/camera_info}. Quindi:
\begin{itemize}[leftmargin=1.4em]
  \item finche' \topic{/camera_info} non arriva, non esistono vere misure 3D;
  \item se la calibrazione camera cambiasse a runtime, questo package non la rilegge.
\end{itemize}

\subsubsection{Passo 2: class filtering e costruzione dell'oggetto Armor}

Per ogni detection, \code{DetectionConverter::buildArmorFromDetection()} esegue:
\begin{enumerate}[leftmargin=1.4em]
  \item lettura di centro, larghezza, altezza della bbox;
  \item correzione verticale tramite \code{bbox\_padding\_y};
  \item scarto di bbox degeneri;
  \item filtro per \code{target\_classes};
  \item costruzione di quattro punti 2D sintetici per le due light bar.
\end{enumerate}

Questo e' un punto essenziale: il package non riceve corner precisi delle light bar. Li \emph{approssima} prendendo i lati della bbox e restringendoli con \code{light\_ratio}. In altre parole, il PnP si appoggia a una geometria sintetica, non a keypoint reali.

Nel codice attuale il tipo SMALL/LARGE non viene deciso dall'aspect ratio della bbox. La detection viene prima convertita in un'armatura interna con tipo provvisorio, poi \code{PnPSolver::solveBestArmorType()} prova entrambi i modelli fisici e sceglie quello con errore medio di reproiezione piu' basso. Questo rende la scelta del tipo piu' legata alla geometria camera-modello e meno a un rapporto larghezza/altezza rumoroso.

\subsubsection{Passo 3: PnP}

Il \code{PnPSolver} usa:
\begin{itemize}[leftmargin=1.4em]
  \item armatura piccola: \(140 \times 125\) mm;
  \item armatura grande: \(235 \times 127\) mm;
  \item quattro punti co-planari ordinati in senso orario;
  \item \code{cv::SOLVEPNP\_IPPE}, adatto a oggetti planari.
\end{itemize}

Il risultato e' una posa 3D nel frame camera:
\begin{equation}
  \mathbf{t} = [x, y, z]^T, \qquad R = \text{Rodrigues}(\mathbf{r})
\end{equation}

Il package converte poi la rotazione in quaternion ROS e salva anche la distanza dal centro immagine.

Dopo la solve, il package riproietta i quattro punti 3D nel piano immagine e calcola l'errore medio in pixel. Se questo errore supera \code{pnp.max\_reprojection\_error}, la detection viene scartata. Per un principiante questo e' un controllo da ricordare: non basta che OpenCV produca una posa; la posa deve anche spiegare bene i punti 2D osservati.

\subsubsection{Passo 4: trasformazione al frame di tracking}

Una volta ottenuta la posa in frame camera, il package fa un lookup TF:
\begin{equation}
  T_{\text{target\_frame} \leftarrow \text{camera}}
\end{equation}

e trasforma ogni armatura nel frame di tracking. Se il TF fallisce, l'intera lista di armature del frame viene svuotata. Questo significa che un problema TF non degrada ``un po''' la pipeline: la interrompe a livello di measurement.

\subsubsection{Passo 5: filtro di plausibilita'}

Il package scarta una detection 3D se:
\begin{enumerate}[leftmargin=1.4em]
  \item \(|z| > 1.2\) m;
  \item la norma planare \(\sqrt{x^2 + y^2}\) supera \code{max\_armor\_distance}.
\end{enumerate}

Questo filtro e' semplice ma molto pratico: taglia pose numericamente assurde prima che entrino nel tracking.

\subsubsection{Figura geometrica del blocco di misura}

\begin{figure}[ht]
\centering
\begin{tikzpicture}[
  scale=1.0,
  arr/.style={-{Latex[length=2.5mm]}, thick},
  lbl/.style={font=\small}
]
  \draw[fill=black!3] (0,0) rectangle (6,4);
  \node at (3,4.3) {\small Immagine detector};
  \draw[thick,blue!70] (1.3,0.8) rectangle (4.7,3.2);
  \node[lbl,blue!70!black] at (3,3.45) {bbox YOLO};

  \draw[thick,red!70] (1.7,1.1) -- (1.7,2.9);
  \draw[thick,red!70] (4.3,1.1) -- (4.3,2.9);
  \fill[red!70] (1.7,1.1) circle (1.4pt);
  \fill[red!70] (1.7,2.9) circle (1.4pt);
  \fill[red!70] (4.3,2.9) circle (1.4pt);
  \fill[red!70] (4.3,1.1) circle (1.4pt);
  \node[lbl,red!70!black] at (3,0.45) {4 punti sintetici per PnP};

  \draw[arr] (6.4,2) -- (8.3,2);
  \node[lbl] at (7.35,2.35) {solvePnP};

  \draw[fill=green!8] (8.7,0.8) rectangle (13.2,3.2);
  \coordinate (camorigin) at (10.0,1.7);
  \draw[arr] (camorigin) -- ++(2.0,0) node[right] {\small \(z\) ottico};
  \draw[arr] (camorigin) -- ++(0,1.2) node[above] {\small \(-y\) ottico};
  \draw[arr] (camorigin) -- ++(-1.1,-0.8) node[below left] {\small \(x\) ottico};
  \draw[thick] (11.2,1.7) ellipse (0.9 and 0.5);
  \node[lbl] at (11.2,0.9) {posa 3D in frame camera};
\end{tikzpicture}
\caption{Dal bounding box alla posa 3D: il package approssima le light bar e poi risolve un PnP planare.}
\end{figure}

\subsection{Sottosistema posa e TF: come il package sa dove si trova la camera}

\subsubsection{Tre modalita' di sorgente posa}

Il package dichiara tre modalita' operative per \code{pose\_source}:
\begin{enumerate}[leftmargin=1.4em]
  \item \code{none};
  \item \code{micro\_pose};
  \item \code{camera\_imu}.
\end{enumerate}

Il comportamento non e' equivalente.

\paragraph{Modalita' \code{micro\_pose}.}
Il callback legge:
\begin{itemize}[leftmargin=1.4em]
  \item pitch da \code{msg->pose.position.x};
  \item yaw da \code{msg->pose.position.y}.
\end{itemize}

Questo e' un dettaglio di integrazione molto importante: il messaggio \code{PoseStamped} non viene interpretato in modo ``semantico ROS classico''. Non si usa il quaternion della pose, ma due campi posizione come contenitori di angoli.

\paragraph{Modalita' \code{camera\_imu}.}
Il callback estrae roll, pitch e yaw direttamente dal quaternion dell'IMU.

\paragraph{Modalita' \code{none}.}
Il package considera la posa sempre \emph{fresh} dal punto di vista logico, ma non riceve aggiornamenti reali di orientazione. La trasformazione camera viene quindi costruita usando l'orientazione interna corrente, che in assenza di input rimane l'identita'. Questa modalita' e' utilissima per debug software e bagfile, ma su hardware reale sarebbe una semplificazione pericolosa se la camera non fosse effettivamente allineata al frame di tracking.

\subsubsection{Freshness, timeout e comportamento in caso di dati stantii}

Il package usa due soglie diverse:
\begin{itemize}[leftmargin=1.4em]
  \item \code{self\_orientation\_timeout}: soglia di warning/throttling nel logger;
  \item \code{pose\_timeout}: soglia logica per \code{pose\_fresh}.
\end{itemize}

Questo significa che il sistema puo' continuare a pubblicare TF usando l'ultima posa nota, ma contemporaneamente dichiarare la posa \emph{stale} verso il fire gate.

Questa scelta e' intelligente:
\begin{enumerate}[leftmargin=1.4em]
  \item non distrugge la geometria interna al primo dropout;
  \item permette ancora di ragionare sul target;
  \item ma impedisce di sparare quando il riferimento del gimbal non e' affidabile.
\end{enumerate}

\subsubsection{Broadcast del TF camera}

Il \code{GimbalPoseAdapter} costruisce il TF camera con:
\begin{itemize}[leftmargin=1.4em]
  \item altezza del pivot \code{gimbal.height};
  \item offset camera \code{camera\_offset\_x} e \code{camera\_offset\_z};
  \item rotazione di convenzione ottica \((-\pi/2, 0, -\pi/2)\).
\end{itemize}

La traslazione finale e':
\begin{equation}
  \mathbf{t}_{camera} = \mathbf{t}_{pivot} + R_{gimbal} \, \mathbf{offset}_{camera}
\end{equation}

Il package inserisce questo TF sia nel broadcaster sia direttamente nel \code{tf2\_buffer}. Quindi il suo uso di TF non dipende interamente da un broadcaster esterno.

\begin{figure}[ht]
\centering
\begin{tikzpicture}[
  scale=1.0,
  arr/.style={-{Latex[length=2.5mm]}, thick},
  lbl/.style={font=\small}
]
  \coordinate (O) at (0,0);
  \draw[arr] (O) -- (5.4,0) node[right] {\small asse \(x\) in \code{target\_frame}};
  \draw[arr] (O) -- (0,4.6) node[above] {\small asse \(z\) in \code{target\_frame}};

  \coordinate (P) at (1.2,1.7);
  \fill (P) circle (1.5pt);
  \node[lbl, left] at (P) {pivot gimbal};
  \draw[dashed] (P) -- (1.2,0) node[below] {\small \code{gimbal.height}};

  \draw[thick,blue!70] (P) -- ++(1.7,0.8);
  \node[lbl,blue!70!black] at (2.55,2.7) {\small orientazione gimbal};

  \coordinate (C) at ($(P)+(1.0,0.55)$);
  \fill[red!70] (C) circle (1.6pt);
  \node[lbl, right] at (C) {camera};
  \draw[thick,red!70,->] (P) -- (C);
  \node[lbl] at (2.0,1.75) {\small offset camera ruotato};

  \coordinate (T) at (4.2,2.9);
  \fill[green!60!black] (T) circle (1.7pt);
  \node[lbl, right] at (T) {target in \code{odom}};
  \draw[arr,green!50!black] (C) -- (T);
  \node[lbl] at (3.5,3.35) {\small misura nel frame camera};
\end{tikzpicture}
\caption{Schema semplificato della costruzione del TF camera a partire da pivot, offset e orientazione del gimbal.}
\end{figure}

\subsection{Sottosistema di tracking: il cuore matematico del package}

\subsubsection{Perche' il tracker non stima la piastra visibile ma il centro del robot}

Se si stimasse solo la piastra visibile, appena il robot ruota di \(90^\circ\) si perderebbe consistenza: la faccia cambia, il punto osservato cambia, la dinamica apparente salta.

Il package adotta invece un modello ``spinning top'' del robot:
\begin{itemize}[leftmargin=1.4em]
  \item esiste un centro del robot;
  \item le armature stanno su un'orbita attorno al centro;
  \item il robot puo' traslare e ruotare;
  \item la faccia osservata e' una manifestazione di quello stato, non lo stato stesso.
\end{itemize}

\subsubsection{Stato dell'EKF}

Lo stato principale del tracker e':
\begin{equation}
  \mathbf{x} =
  \begin{bmatrix}
    x_c & v_{x_c} & y_c & v_{y_c} & z_a & v_{z_a} & \psi & \dot{\psi} & r
  \end{bmatrix}^{T}
\end{equation}

dove:
\begin{itemize}[leftmargin=1.4em]
  \item \(x_c, y_c\): centro del robot sul piano;
  \item \(v_{x_c}, v_{y_c}\): velocita' del centro;
  \item \(z_a\): quota dell'armatura attualmente modellata;
  \item \(v_{z_a}\): velocita' verticale;
  \item \(\psi\): yaw del corpo;
  \item \(\dot{\psi}\): velocita' angolare di yaw;
  \item \(r\): raggio dell'armatura attiva rispetto al centro.
\end{itemize}

Il modello di osservazione usato dall'EKF e':
\begin{equation}
  h(\mathbf{x}) =
  \begin{bmatrix}
    x_c - r \cos \psi \\
    y_c - r \sin \psi \\
    z_a \\
    \psi
  \end{bmatrix}
\end{equation}

Questa formula e' la chiave di lettura di tutto il tracker: la misura osservata non e' il centro, ma la faccia del robot prodotta da centro + yaw + raggio.

\subsubsection{Dinamica e damping}

La factory dei tracker costruisce l'EKF iniettando:
\begin{itemize}[leftmargin=1.4em]
  \item modello di processo \(f\);
  \item Jacobiano \(J_f\);
  \item modello di misura \(h\);
  \item Jacobiano \(J_h\);
  \item rumore di processo \(Q\);
  \item rumore di misura \(R\).
\end{itemize}

Il modello di processo non e' un semplice constant-velocity puro. Applica un \emph{velocity decay} separato per:
\begin{itemize}[leftmargin=1.4em]
  \item velocita' traslazionali;
  \item velocita' di yaw.
\end{itemize}

Questo damping viene adattato frame per frame da \code{computeAdaptiveDamping()} in base a innovazione, overshoot, stato TEMP\_LOST e stationary detection. In altre parole: il package non si limita a stimare uno stato, ma modula quanto crede alla propria dinamica quando sente che sta sovra-predicendo.

\subsubsection{Rumore di misura dipendente da distanza e obliquita'}

La matrice \(R\) non e' costante. Dipende da:
\begin{enumerate}[leftmargin=1.4em]
  \item distanza del target;
  \item obliquita' della faccia rispetto alla camera.
\end{enumerate}

Questo e' corretto fisicamente: una faccia lontana o molto obliqua produce un PnP meno affidabile. Il package aumenta quindi il rumore percepito della misura quando la geometria osservativa peggiora.

Quando l'obliquita' supera \code{max\_yaw\_oblique\_deg}, il rumore sullo yaw viene praticamente inflazionato a un valore enorme. Conseguenza pratica: la misura di yaw viene quasi ignorata, pur continuando a fondere la posizione.

\subsubsection{Associazione detection-tracker}

Il tracker manager usa il seguente schema:
\begin{enumerate}[leftmargin=1.4em]
  \item predice tutti i tracker;
  \item per ogni detection cerca il tracker con costo di assegnazione minimo;
  \item il costo e' la distanza di Mahalanobis;
  \item tracker in stato DETECTING vengono penalizzati con un fattore \(4\).
\end{enumerate}

La distanza di Mahalanobis e':
\begin{equation}
  d^2 = (\mathbf{z} - h(\mathbf{x}_{pri}))^T S^{-1} (\mathbf{z} - h(\mathbf{x}_{pri}))
\end{equation}
con
\begin{equation}
  S = H P_{pri} H^T + R
\end{equation}

Questo e' molto piu' informativo di una distanza euclidea pura, perche' tiene conto dell'incertezza anisotropa del filtro.

\subsubsection{Geometria avanzata: due raggi e offset verticale}

Per i robot a quattro facce, il package modella due coppie di facce con:
\begin{itemize}[leftmargin=1.4em]
  \item \code{radius\_1}: raggio della coppia attiva;
  \item \code{radius\_2}: raggio dell'altra coppia;
  \item \code{dz}: offset verticale tra coppie.
\end{itemize}

Ma questi valori non sono tutti dentro l'EKF principale. Il package tiene:
\begin{itemize}[leftmargin=1.4em]
  \item un EKF 9D per lo stato principale;
  \item due scalar KF esterni per i due raggi;
  \item una EMA per \code{dz}.
\end{itemize}

Questa scelta e' sofisticata ma sensata:
\begin{enumerate}[leftmargin=1.4em]
  \item evita di gonfiare troppo la dimensione dell'EKF;
  \item permette gain adattivo sui raggi, che cambiano lentamente;
  \item rende piu' robusto il cambio di coppia durante i face jump.
\end{enumerate}

\subsubsection{Face jump}

Un punto davvero avanzato del package e' la gestione del \emph{face jump}: quando la posizione osservata e' plausibile, ma lo yaw salta abbastanza da indicare che non stiamo piu' vedendo la stessa faccia.

Quando questo succede:
\begin{enumerate}[leftmargin=1.4em]
  \item il package ``snap''-pa lo yaw allo yaw misurato;
  \item se il salto e' circa \(90^\circ\), interpreta un cambio di coppia;
  \item scambia i due scalar KF dei raggi;
  \item aggiorna \code{dz};
  \item se il centro implicito diverge troppo dalla detection, effettua un hard reset dello stato;
  \item poi esegue comunque un update EKF con la nuova misura.
\end{enumerate}

Questo e' il motivo per cui il tracker resta coerente anche su robot che ruotano rapidamente.

\subsubsection{Secondary face fusion}

Se due facce dello stesso robot sono visibili nello stesso frame, il package prova a usare anche la faccia secondaria con un update sequenziale. In pratica:
\begin{enumerate}[leftmargin=1.4em]
  \item usa la faccia primaria per l'update principale;
  \item costruisce un modello di osservazione alternativo per la faccia secondaria;
  \item applica un secondo update EKF, se il Mahalanobis della faccia secondaria e' accettabile.
\end{enumerate}

Questa e' una tecnica molto avanzata: sfrutta informazione geometrica ridondante senza introdurre un enorme modello multi-osservazione unico.

\subsubsection{Stationary detection}

Il tracker misura anche se il target appare sostanzialmente fermo, combinando:
\begin{itemize}[leftmargin=1.4em]
  \item delta di misura tra frame;
  \item norma dell'innovazione;
  \item velocita' lineare stimata;
  \item velocita' angolare stimata.
\end{itemize}

Quando la confidenza di stazionarieta' cresce, il package collassa progressivamente le velocita' verso zero. Questo riduce il rischio di ``coast infinito'' su target fermi.

\subsubsection{Macchina a stati del tracker}

\begin{figure}[ht]
\centering
\begin{tikzpicture}[
  node distance=18mm and 16mm,
  state/.style={draw, rounded corners=2pt, minimum width=2.8cm, minimum height=0.95cm, align=center, fill=blue!6},
  arr/.style={-{Latex[length=2.5mm]}, thick},
  lbl/.style={font=\small, fill=white, inner sep=1pt}
]
  \node[state] (lost) {LOST};
  \node[state, right=of lost] (detecting) {DETECTING};
  \node[state, right=of detecting] (tracking) {TRACKING};
  \node[state, below=of tracking] (temp) {TEMP\_LOST};

  \draw[arr] (lost) -- node[lbl, above] {spawn / init} (detecting);
  \draw[arr] (detecting) -- node[lbl, above] {enough consecutive matches} (tracking);
  \draw[arr] (detecting) -- node[lbl, below] {miss} (lost);
  \draw[arr] (tracking) -- node[lbl, right] {miss} (temp);
  \draw[arr] (temp) -- node[lbl, left] {match returns} (tracking);
  \draw[arr] (temp) -- node[lbl, below] {too many misses} (lost);
\end{tikzpicture}
\caption{Macchina a stati del tracker. La scelta e' conservativa: si conferma lentamente, ma si perde lentamente solo dopo essere stati davvero in tracking.}
\end{figure}

Conseguenza pratica fondamentale:
\begin{itemize}[leftmargin=1.4em]
  \item in DETECTING il target esiste internamente ma non e' ancora ``shootable'';
  \item in TEMP\_LOST il target resta pubblicato come tracking valido, ma il fire gate impedisce il fuoco.
\end{itemize}

\subsubsection{Limitazione strutturale attuale: numero di armature}

Il codice contiene supporto concettuale per target a 2, 3 e 4 facce, ma l'attuale implementazione del tracker imposta sempre:
\begin{equation}
  \texttt{tracked\_armors\_num} = \texttt{NORMAL\_4}
\end{equation}

Quindi il planner ha gia' le astrazioni per altri casi, ma il runtime attuale non sta realmente inferendo il tipo robot a partire dalla classificazione o da un'altra sorgente.

\begin{figure}[ht]
\centering
\begin{tikzpicture}[
  scale=1.0,
  arr/.style={-{Latex[length=2.5mm]}, thick},
  lbl/.style={font=\small}
]
  \coordinate (C) at (0,0);
  \draw[fill=blue!8] (C) circle (1.35);
  \fill[black] (C) circle (1.5pt);
  \node[lbl] at (0,-1.7) {centro del robot};

  \coordinate (A0) at (1.0,0.0);
  \coordinate (A1) at (0.0,1.25);
  \coordinate (A2) at (-1.0,0.0);
  \coordinate (A3) at (0.0,-1.25);

  \draw[thick,green!60!black] (C) -- (A0);
  \draw[thick,orange!80!black] (C) -- (A1);
  \draw[thick,green!60!black] (C) -- (A2);
  \draw[thick,orange!80!black] (C) -- (A3);

  \draw[fill=green!25] ($(A0)+(0.16,0)$) rectangle ($(A0)+(-0.16,0.28)$);
  \draw[fill=orange!25] ($(A1)+(0.28,0.16)$) rectangle ($(A1)+(0,-0.16)$);
  \draw[fill=green!25] ($(A2)+(-0.16,0)$) rectangle ($(A2)+(0.16,-0.28)$);
  \draw[fill=orange!25] ($(A3)+(-0.28,-0.16)$) rectangle ($(A3)+(0,0.16)$);

  \node[lbl,green!50!black] at (1.65,0.4) {\small coppia A: \(r_1\)};
  \node[lbl,orange!70!black] at (1.9,1.35) {\small coppia B: \(r_2\), \(dz\)};

  \draw[arr] (0,0) -- (0.75,0.55);
  \node[lbl] at (0.9,0.8) {\small yaw corpo \(\psi\)};
\end{tikzpicture}
\caption{Modello ``spinning top'': il tracker stima il centro e la geometria delle facce, non una singola piastra indipendente.}
\end{figure}

\subsection{Sottosistema planning: prevedere, scegliere, decidere}

\subsubsection{Latenza: due usi diversi}

Il \code{LatencyCompensator} fa due cose diverse, che e' importante non confondere:
\begin{enumerate}[leftmargin=1.4em]
  \item aggiorna una stima smussata di \code{time\_bias};
  \item calcola il \code{transportDelay} grezzo del frame.
\end{enumerate}

Nel package queste due quantita' vengono usate in modo diverso:
\begin{itemize}[leftmargin=1.4em]
  \item \code{time\_bias} entra nella previsione in avanti della soluzione di tiro;
  \item \code{transportDelay} entra nell'eta' effettiva della misura e quindi nella sua staleness.
\end{itemize}

Questo e' un dettaglio molto raffinato: la stessa latenza non viene usata in modo monolitico, ma separata in componente predittiva e componente diagnostica.

\subsubsection{Ballistica}

Il \code{BallisticsSolver} risolve iterativamente pitch e flight time. Non e' un simulatore balistico completo a 6 gradi di liberta'; usa due modelli semplificati:
\begin{enumerate}[leftmargin=1.4em]
  \item drag lineare;
  \item drag quadratico ispirato alla perdita di velocita' sull'orizzontale.
\end{enumerate}

Il cuore dell'iterazione e':
\begin{align}
  t_0 &\approx \frac{d}{v_{muzzle}} \\
  \theta &\approx \arctan\left(\frac{z_{target} + \frac{1}{2} g t^2}{d}\right)
\end{align}

poi il solver aggiorna \(t\) con una formula che dipende dal modello di drag scelto.

Questa distinzione e' importante:
\begin{itemize}[leftmargin=1.4em]
  \item \code{valid} significa che il modello ha prodotto numeri finiti e coerenti;
  \item \code{reachable} significa che il pitch richiesto rientra nei limiti meccanici del gimbal.
\end{itemize}

\subsubsection{Tre modalita' di tiro}

Il planner puo' costruire tre tipi di piano:
\begin{enumerate}[leftmargin=1.4em]
  \item \textbf{VISIBLE\_DIRECT}: mira la faccia attualmente visibile, se e' fresca e osservabile;
  \item \textbf{PREDICTED\_DIRECT}: se la faccia corrente non e' piu' affidabile, predice quale faccia sara' migliore a breve;
  \item \textbf{INDIRECT}: non aspetta di vedere bene una faccia, ma sincronizza il tiro con l'istante in cui una faccia in rotazione entrera' nella finestra utile.
\end{enumerate}

\begin{figure}[ht]
\centering
\begin{tikzpicture}[
  node distance=8mm and 12mm,
  box/.style={draw, rounded corners=2pt, minimum width=3.0cm, minimum height=0.95cm, align=center, fill=blue!6},
  choice/.style={draw, diamond, aspect=2.0, align=center, inner sep=1pt, fill=orange!10},
  arr/.style={-{Latex[length=2.5mm]}, thick},
  lbl/.style={font=\small, fill=white, inner sep=1pt}
]
  \node[box] (start) {TrackSnapshot};
  \node[choice, below=of start] (spin) {spin alto?};
  \node[choice, below left=of spin] (visible) {faccia matched\\fresca e osservabile?};
  \node[box, below left=of visible] (vd) {VISIBLE\_DIRECT};
  \node[box, below right=of visible] (pd) {PREDICTED\_DIRECT};
  \node[box, below right=of spin] (ind) {INDIRECT};
  \node[box, below=18mm of spin] (score) {score multi-obiettivo};
  \node[box, below=of score] (gate) {FireGate};

  \draw[arr] (start) -- (spin);
  \draw[arr] (spin) -- node[lbl, left] {no} (visible);
  \draw[arr] (spin) -- node[lbl, right] {si} (ind);
  \draw[arr] (visible) -- node[lbl, left] {si} (vd);
  \draw[arr] (visible) -- node[lbl, right] {no} (pd);
  \draw[arr] (vd) |- (score);
  \draw[arr] (pd) |- (score);
  \draw[arr] (ind) |- (score);
  \draw[arr] (score) -- (gate);
\end{tikzpicture}
\caption{Logica ad alto livello delle modalita' di planning.}
\end{figure}

\subsubsection{Visible direct}

Questa modalita' usa la \code{matched\_face} del tracker. Quindi e' la forma di ingaggio piu' aderente all'osservazione corrente. Viene scelta solo se la faccia e':
\begin{itemize}[leftmargin=1.4em]
  \item valida;
  \item fresca;
  \item osservabile.
\end{itemize}

In pratica e' la modalita' migliore quando la misura visuale e' ancora affidabile e il target non gira troppo.

\subsubsection{Predicted direct}

Qui il planner esamina tutte le facce teoriche del target, ne predice la posizione futura a un certo orizzonte, e costruisce un piano per ciascuna. Poi sceglie la migliore.

Questo serve quando la faccia attualmente matched non e' piu' abbastanza affidabile, ma il modello dinamico del target e' ancora credibile.

\subsubsection{Indirect}

Questa e' la modalita' piu' sofisticata. Si usa quando la velocita' angolare del target supera una soglia. Il planner non cerca piu' la faccia migliore ``ora'', ma il prossimo istante di allineamento utile:
\begin{equation}
  t_{align} = \text{nextAlignmentTime}(\psi, \dot{\psi}, \text{face\_index}, \Delta\psi, \text{bearing})
\end{equation}

Poi confronta \(t_{align}\) con:
\begin{equation}
  t_{desired} = t_{flight} + t_{bias} + t_{gimbal\_response}
\end{equation}

e minimizza il residuo temporale.

Questo permette di sparare ``nel futuro'' verso una faccia che non e' quella otticamente migliore in questo istante, ma che diventera' corretta all'arrivo del proiettile.

\subsubsection{Engagement planner e funzione di costo}

L'\code{EngagementPlanner} prende il piano generato per ogni target e ne calcola un punteggio. Il costo include:
\begin{itemize}[leftmargin=1.4em]
  \item distanza;
  \item tempo di volo;
  \item incertezza angolare;
  \item costo di slew del gimbal;
  \item staleness della misura;
  \item bassa visibilita';
  \item penalita' per temp\_lost;
  \item penalita' per switch di target;
  \item penalita' per switch di faccia sullo stesso target;
  \item penalita' per margine di fuoco negativo;
  \item penalita' per piano non raggiungibile.
\end{itemize}

I piani senza balistica valida o fuori dal range minimo/massimo vengono scartati prima dello scoring. Questo evita di confrontare piani che non dovrebbero mai arrivare al fire gate.

In forma schematica:
\begin{equation}
  \text{score} =
  w_r C_{range} +
  w_f C_{flight} +
  w_u C_{uncertainty} +
  w_s C_{slew} +
  w_{st} C_{staleness} +
  w_v C_{visibility} + \dots
\end{equation}

Il target finale e' quello con \emph{score minimo}.

\paragraph{Dettaglio avanzato importante.}
Il nodo attuale dichiara anche i parametri ROS \code{cost.*}: \code{cost.range}, \code{cost.flight\_time}, \code{cost.uncertainty}, \code{cost.slew}, \code{cost.switch\_target}, \code{cost.switch\_face}, \code{cost.staleness}, \code{cost.temp\_lost}, \code{cost.low\_visibility} e \code{cost.negative\_margin}. Quindi i pesi della cost function sono modificabili dal launch, ma vanno toccati dopo aver validato measurement, TF, tracking e latenza.

\subsubsection{Fire gate}

Anche un buon piano non implica automaticamente il fuoco. Il \code{FireGate} controlla in ordine:
\begin{enumerate}[leftmargin=1.4em]
  \item validita' dello stato di tracking;
  \item stato TEMP\_LOST;
  \item staleness della misura;
  \item validita' balistica;
  \item raggiungibilita' del pitch;
  \item range minimo/massimo;
  \item margine della fire window;
  \item freschezza della posa gimbal, se una pose source e' configurata;
  \item consenso esplicito \code{debug.allow\_fire\_without\_pose}, se \code{pose\_source=none}.
\end{enumerate}

Il risultato non e' solo booleano: contiene anche un \code{blocker} testuale, molto utile in debug.

\subsection{Output: come il package materializza la decisione}

\subsubsection{Target vs Targets}

Il package pubblica sia:
\begin{itemize}[leftmargin=1.4em]
  \item un messaggio \code{Targets} con tutti i target attivi;
  \item un messaggio \code{Target} singolo per il best tracker.
\end{itemize}

Questo e' comodo perche':
\begin{enumerate}[leftmargin=1.4em]
  \item un consumer semplice puo' leggere solo il best target;
  \item un consumer avanzato puo' ragionare sull'insieme completo.
\end{enumerate}

\subsubsection{Comandi gimbal}

Il \code{GimbalCommandController} fa tre cose:
\begin{enumerate}[leftmargin=1.4em]
  \item clamp degli angoli assoluti pitch/yaw;
  \item conversione del comando principale in gradi;
  \item costruzione di un marker di impatto.
\end{enumerate}

Dettaglio molto importante: il comando pubblicato su \topic{/tracker/cmd_gimbal} usa angoli \emph{assoluti} nel frame di riferimento, convertiti in \emph{gradi}. Il controller downstream conosce la posa corrente del gimbal e calcola la correzione necessaria. Questa scelta evita di dipendere da feedback gimbal perfettamente fresco per costruire comandi relativi.

Ancora piu' importante dal punto di vista di integrazione: il \code{Twist} pubblicato su \topic{/cmd_vel} non rappresenta una twist standard. Nel codice corrente e' un canale legacy parziale:
\begin{itemize}[leftmargin=1.4em]
  \item \code{angular.y} contiene il pitch assoluto clampato in radianti;
  \item \code{angular.z} contiene lo yaw assoluto clampato in radianti;
  \item il flag di fuoco affidabile e' \code{GimbalCmd::fire\_cmd}, non \code{Twist::angular.x}.
\end{itemize}

Questa scelta e' utile per compatibilita', ma va capita bene per non interpretare male il topic.

\subsubsection{Perche' non c'e' smoothing nel command controller}

Il codice corrente non applica EMA o altri smoothing nel \code{GimbalCommandController}. I commenti del file spiegano la scelta: la convergenza viene lasciata al PID o al controller downstream del gimbal. Questo mantiene il nodo di targeting piu' diretto: decide \emph{dove} puntare e \emph{se} sparare, ma non cerca di filtrare dinamicamente il movimento dell'attuatore.

\subsubsection{Debug e RViz}

Il package costruisce:
\begin{itemize}[leftmargin=1.4em]
  \item marker del centro del target;
  \item freccia di velocita' lineare;
  \item freccia di velocita' angolare;
  \item marker per le armature stimate;
  \item marker dell'impatto previsto.
\end{itemize}

Questi marker sono estremamente utili per capire se l'errore sta nella measurement, nel tracking o nel planner.

\subsection{Parametri importanti: default del nodo vs override del launch}

Il nodo dichiara i parametri in \path{src/auto_aim_targeting/src/auto_aim_node.cpp}. Un launch puo' poi sovrascriverli. La tabella seguente riassume alcuni parametri ad alto impatto.

\begin{longtable}{p{0.24\textwidth} p{0.13\textwidth} p{0.14\textwidth} p{0.39\textwidth}}
\toprule
\textbf{Parametro} & \textbf{Default nodo} & \textbf{debug\_targeting} & \textbf{Significato tecnico} \\
\midrule
\endfirsthead
\toprule
\textbf{Parametro} & \textbf{Default nodo} & \textbf{debug\_targeting} & \textbf{Significato tecnico} \\
\midrule
\endhead
\code{target\_frame} & \code{odom} & \code{odom} & Frame in cui vengono trasformate e tracciate le armature. \\
\code{target\_classes} & \code{\{"3"\}} & \code{\{"3"\}} & Insieme di classi detector accettate come target. \\
\code{light\_ratio} & 0.85 & 1.0 & Quanto la bbox viene ``stretta'' per ricostruire i punti PnP sintetici. \\
\code{bbox\_padding\_y} & 80.0 & non ridefinito & Correzione verticale della bbox, tipicamente legata a letterboxing. \\
\code{pnp.max\_reprojection\_error} & 10.0 px & non ridefinito & Errore medio massimo accettato dopo la riproiezione dei punti PnP. \\
\code{max\_armor\_distance} & 10.0 & 10.0 & Limite massimo di distanza per accettare una detection 3D. \\
\code{tracker.max\_match\_distance} & 0.15 & 0.30 & Gate metrico grezzo per match e jump handling. \\
\code{tracker.tracking\_thres} & 5 & 4 & Numero di match consecutivi richiesti per promuovere DETECTING -> TRACKING. \\
\code{tracker.lost\_time\_thres} & 0.30 s & 0.50 s & Tempo massimo di coast in TEMP\_LOST. \\
\code{tracker.max\_trackers} & 5 & 5 & Numero massimo di tracker simultanei. \\
\code{ekf.sigma2\_q\_xyz} & 5.0 & 1.0 & Rumore di processo traslazionale. \\
\code{ekf.sigma2\_q\_yaw} & 10.0 & 3.0 & Rumore di processo angolare. \\
\code{ekf.xyz\_damping\_alpha} & 0.95 & 0.85 & Decadimento base delle velocita' lineari. \\
\code{ekf.yaw\_damping\_alpha} & 0.95 & 0.85 & Decadimento base della velocita' di yaw. \\
\code{pose\_source} & \code{none} & \code{none} & Sorgente di posa del gimbal/camera. \\
\code{micro\_pose\_timeout} & 0.15 & non ridefinito & Soglia di freshness della posa per il fire gate. \\
\code{debug.allow\_fire\_without\_pose} & \code{false} & \code{true} & Permette il fuoco anche senza pose source configurata; utile in test, da trattare con cautela su hardware reale. \\
\code{gimbal.height} & 0.325 & 0.325 & Altezza del pivot del gimbal nel frame di tracking. \\
\code{camera\_offset\_x} & 0.107 & 0.107 & Offset camera lungo l'asse del supporto. \\
\code{camera\_offset\_z} & 0.136 & 0.136 & Offset camera verticale rispetto al pivot. \\
\code{bullet\_speed} & 25.0 & 25.0 & Velocita' iniziale del proiettile. \\
\code{time\_bias} & 0.08 & 0.05 & Offset predittivo di latenza. \\
\code{time\_bias\_alpha} & 0.35 & 0.35 & Fattore EMA per aggiornare la latenza stimata. \\
\code{angular\_window} & 0.09 & 0.09 & Finestra angolare base per la fire logic. \\
\code{indirect\_vyaw\_threshold} & 3.0 & 3.0 & Soglia di spin oltre cui conviene considerare tiro indiretto. \\
\code{indirect\_timing\_tolerance} & 0.02 & 0.02 & Errore temporale tollerato in indirect mode. \\
\code{cost.switch\_face} & 0.20 & 0.20 & Penalita' per cambiare faccia sullo stesso tracker, utile contro jitter tra facce candidate. \\
\code{max\_cmd\_angle} & 30.0 deg & 30.0 deg & Clamp del pitch assoluto pubblicato. \\
\code{max\_cmd\_yaw\_angle} & 180.0 deg & 180.0 deg & Clamp dello yaw assoluto pubblicato. \\
\bottomrule
\end{longtable}

\paragraph{Regola pratica di tuning.}
L'ordine corretto e':
\begin{enumerate}[leftmargin=1.4em]
  \item measurement;
  \item pose/TF;
  \item tracking;
  \item latenza;
  \item planner, fire gate e indirect mode;
  \item clamp e convenzioni di output dei comandi.
\end{enumerate}

Fare l'opposto porta quasi sempre a compensare con parametri a valle errori nati a monte.

\subsection{Ragionamento critico: perche' le scelte del package sono sensate}

\subsubsection{Perche' EKF + scalar KF invece di un unico EKF gigante}

Una domanda da livello avanzato e' se non sarebbe piu' ``pulito'' mettere tutto in uno stato unico: centro, velocita', yaw, due raggi, offset verticale, eccetera.

In teoria si puo'. In pratica questo package sceglie una via piu' robusta:
\begin{itemize}[leftmargin=1.4em]
  \item EKF principale per dinamica del centro e dello yaw;
  \item scalar KF per parametri geometrici lenti;
  \item EMA per offset verticale osservato meno spesso.
\end{itemize}

Questa decomposizione riduce complessita' numerica, facilita il face jump e rende piu' semplice il reset selettivo dei soli sotto-modelli che stanno divergendo.

\subsubsection{Perche' Mahalanobis e non distanza euclidea}

La distanza euclidea sa solo ``quanto e' lontano'' un punto. La Mahalanobis sa anche ``quanto sarebbe improbabile vederlo li' dato lo stato e l'incertezza attuale del filtro''.

Per un tracker probabilistico questa e' la metrica giusta.

\subsubsection{Perche' bypassare il gate di yaw a forte obliquita'}

Quando una faccia e' quasi edge-on, lo yaw ottenuto da PnP diventa poco informativo. Se si insistesse ad applicare lo yaw gate in modo rigido, il tracker genererebbe falsi face jump. Il package quindi:
\begin{enumerate}[leftmargin=1.4em]
  \item rileva l'obliquita';
  \item evita di classificare come jump un semplice peggioramento geometrico dell'osservazione;
  \item aumenta il rumore di yaw invece di fidarsi di un angolo cattivo.
\end{enumerate}

Questa e' una scelta molto matura dal punto di vista probabilistico.

\subsubsection{Perche' il fire gate e' severo anche quando il planner ha un piano}

Il planner sceglie il \emph{meno peggio} tra i piani disponibili. Ma il meno peggio non e' automaticamente un piano sparabile.

Per questo il fire gate e' separato: serve come filtro finale di sicurezza e consistenza.

\section{Approfondimenti tecnici aggiuntivi}

\subsection{Anatomia dei messaggi principali}

Un altro modo molto utile di capire \pkg\ e' leggere la semantica dei suoi messaggi pubblicati. Il package comunica infatti uno stato molto piu' ricco di un semplice ``bersaglio in \(x,y,z\)''.

\paragraph{Messaggio \code{Target}.}
Il messaggio \code{auto\_aim\_interfaces::msg::Target} contiene:
\begin{itemize}[leftmargin=1.4em]
  \item metadati del tracker: \code{tracking}, \code{id}, \code{tracker\_id}, \code{tracker\_state}, \code{armors\_num};
  \item stato cinematico: posizione, velocita', accelerazione;
  \item stato angolare: \code{yaw}, \code{v\_yaw};
  \item geometria del bersaglio: \code{radius\_1}, \code{radius\_2}, \code{dz};
  \item incertezze: \code{v\_yaw\_variance}, \code{position\_variance\_x/y/z};
  \item ultimo istante in cui e' stata fusa una misura vera.
\end{itemize}

Questo significa che un consumer a valle non riceve solo ``dove puntare'', ma anche \emph{quanto fidarsi} della stima.

\begin{lstlisting}[style=shellstyle,caption={Estratto semantico del messaggio Target}]
bool tracking
string id
int32 tracker_id
int32 armors_num
uint8 tracker_state

float64 v_yaw_variance
float64 position_variance_x
float64 position_variance_y
float64 position_variance_z

geometry_msgs/Point position
geometry_msgs/Vector3 velocity
geometry_msgs/Vector3 acceleration
float64 yaw
float64 v_yaw
float64 radius_1
float64 radius_2
float64 dz
builtin_interfaces/Time last_measurement_stamp
\end{lstlisting}

\paragraph{Messaggio \code{Targets}.}
Questo messaggio e' semplicemente un contenitore con:
\begin{itemize}[leftmargin=1.4em]
  \item array di \code{Target};
  \item \code{best\_target\_idx}.
\end{itemize}

Dal punto di vista architetturale e' una scelta elegante: il nodo non obbliga i consumer a fidarsi della sua selezione finale, ma espone anche il set completo dei candidati.

\paragraph{Messaggio \code{GimbalCmd}.}
Il messaggio \code{GimbalCmd} contiene:
\begin{itemize}[leftmargin=1.4em]
  \item \code{pitch};
  \item \code{yaw};
  \item \code{distance};
  \item \code{fire\_cmd}.
\end{itemize}

Il dettaglio da non dimenticare e' che \code{pitch} e \code{yaw} escono in \emph{gradi} e rappresentano angoli assoluti clampati. Questo e' molto facile da dimenticare quando quasi tutta la pipeline interna usa radianti.

\paragraph{Messaggio \code{TrackerInfo}.}
Serve per diagnosi. Non descrive lo stato completo del target, ma:
\begin{itemize}[leftmargin=1.4em]
  \item differenza posizione tra prediction e misura;
  \item differenza yaw;
  \item \code{face\_angle}, cioe' l'obliquita' della faccia;
  \item fattori di damping effettivi usati nel frame corrente.
\end{itemize}

Questo e' esattamente il tipo di messaggio che un operatore usa per capire \emph{se il filtro sta lottando contro la misura} o \emph{se la misura e' pulita ma il planner e' severo}.

\subsection{Piccoli dettagli tecnici facili da perdere}

\begin{enumerate}[leftmargin=1.4em]
  \item \textbf{\code{dt} non e' libero}: in \code{updateTimestep()} viene clampato tra \(0.01\) s e \(0.10\) s. Questo evita che una variazione di timestamp patologica faccia esplodere la prediction.
  \item \textbf{Il nodo usa sia \code{input.stamp} sia \code{now()}}: il frame timestamp serve per misura e latenza; \code{now()} serve come tempo locale del nodo. La distinzione e' corretta e importante.
  \item \textbf{\code{cam\_info\_sub\_.reset()}}: dopo il primo \topic{/camera_info}, il nodo smette intenzionalmente di ascoltare. E' una decisione di lifetime management, non un incidente.
  \item \textbf{\code{best\_tracker\_id} e \code{previous\_tracker\_id} sono diversi}: il primo e' l'output del frame corrente, il secondo serve per isteresi, continuita' del planner e penalita' di switch.
  \item \textbf{\code{pose\_fresh} e' indipendente dal fatto che il TF venga ancora pubblicato}: il package puo' continuare a costruire geometria usando l'ultima posa nota, ma vietare comunque il fuoco.
  \item \textbf{Il planner e il fire gate sono separati}: una soluzione puo' essere la migliore disponibile e al tempo stesso non essere sparabile.
  \item \textbf{Il tracker non pubblica in DETECTING}: esiste internamente, ma il sistema evita di trasformare un target ancora fragile in un target attivo verso l'esterno.
\end{enumerate}

\subsection{Snippet di codice chiave}

\subsubsection{Pipeline per-frame del nodo}

Il seguente estratto e' il modo piu' compatto di vedere il passaggio da detection a decisione:

\begin{lstlisting}[style=cppstyle,caption={Estratto reale semplificato da AutoAimNode: callback e pipeline per-frame}]
void AutoAimNode::armorsCallback(
  const vision_msgs::msg::Detection2DArray::ConstSharedPtr detection_msg)
{
  if (!isPipelineReady()) {
    return;
  }

  gimbal_pose_adapter_->updateForFrame(detection_msg->header.stamp);
  const GimbalPoseState pose_state = gimbal_pose_adapter_->currentPoseState(now());
  const Observation input = detection_converter_->buildObservation(detection_msg, pose_state);
  processFrame(input);
}

void AutoAimNode::processFrame(const Observation & input)
{
  updateTimestep(input.stamp);
  Trackers::predictAll(trackers_, dt_, config_.tracking.lost_time_thres);
  const auto detection_owners = Trackers::assignDetections(
    trackers_, input.armors, config_.tracking.tracker.mahalanobis_jump_gate);
  Trackers::updateMatched(trackers_, input.armors, detection_owners);
  Trackers::spawnFromUnmatched(...);
  Trackers::pruneLost(trackers_, get_logger());

  const auto snapshots = Trackers::collectSnapshots(trackers_, dt_);
  decision = engagement_planner_->selectBestPlan(snapshots, ctx, input.stamp);
  const bool has_pose_source = config_.pose.pose_source != "none";
  const FireGateResult fire_result =
    fire_gate_->evaluate(decision->plan, input.pose_state.fresh, has_pose_source);
  AimResult output = output_publisher_->buildAimResult(...);
  output_publisher_->publish(output);
}
\end{lstlisting}

Il punto importante qui non e' la sintassi, ma il fatto che il package abbia una pipeline strettamente ordinata. Se l'ordine cambia, cambia il significato probabilistico del sistema.

\subsubsection{Modello di osservazione del tracker}

Questo e' il cuore geometrico del tracking:

\begin{lstlisting}[style=cppstyle,caption={Estratto reale dalla TrackerFactory: modello di osservazione h(x)}]
auto h = [](const Eigen::VectorXd & x) {
  Eigen::VectorXd z(4);
  const double yaw = x(6);
  const double radius = x(8);
  z(0) = x(0) - radius * std::cos(yaw);
  z(1) = x(2) - radius * std::sin(yaw);
  z(2) = x(4);
  z(3) = x(6);
  return z;
};
\end{lstlisting}

Queste quattro righe formalizzano il modello ``centro + orbita della faccia''. Se non si capisce questa lambda, non si capisce il tracker.

\subsubsection{Fire gate finale}

\begin{lstlisting}[style=cppstyle,caption={Estratto reale da FireGate::evaluate()}]
if (!plan.tracking_valid) {
  result.blocker = "not_tracking";
  return result;
}
if (plan.temp_lost) {
  result.blocker = "temp_lost";
  return result;
}
if (plan.measurement_stale) {
  result.blocker = "stale_measurement";
  return result;
}
if (!plan.ballistic_valid) {
  result.blocker = "invalid_ballistics";
  return result;
}
if (!plan.reachable) {
  result.blocker = "unreachable_pitch";
  return result;
}
if (plan.fire_window_margin < 0.0) {
  result.blocker = "negative_window";
  return result;
}
if (has_pose_source && !pose_fresh) {
  result.blocker = "stale_pose";
  return result;
}
if (!has_pose_source && !this->config_.allow_fire_without_pose) {
  result.blocker = "no_pose_source";
  return result;
}
\end{lstlisting}

Questo snippet mostra bene che il package non usa un unico ``confidence score'' globale, ma una catena ordinata di criteri di esclusione.

\subsubsection{Comando gimbal assoluto}

\begin{lstlisting}[style=cppstyle,caption={Estratto reale da GimbalCommandController::makeShotOutput()}]
double cmd_yaw = std::clamp(
  plan.absolute_yaw, -this->max_cmd_yaw_angle_rad_, this->max_cmd_yaw_angle_rad_);
double cmd_pitch = std::clamp(
  plan.absolute_pitch, -this->max_cmd_pitch_angle_rad_, this->max_cmd_pitch_angle_rad_);

output.cmd.yaw = cmd_yaw * 180.0 / M_PI;
output.cmd.pitch = cmd_pitch * 180.0 / M_PI;
output.cmd.distance = plan.range;
output.cmd.fire_cmd = fire;

output.twist.angular.y = cmd_pitch;
output.twist.angular.z = cmd_yaw;
\end{lstlisting}

Questo snippet chiarisce due dettagli che generano spesso confusione: \code{GimbalCmd} esce in gradi, mentre il \code{Twist} legacy conserva pitch e yaw in radianti; inoltre il comando e' assoluto, non relativo alla posa gimbal corrente.

\subsection{Grafici interpretativi derivati direttamente dal codice}

\subsubsection{Inflazione del rumore di yaw in funzione dell'obliquita'}

Nel modello di misura del tracker, il rumore sullo yaw cresce molto quando la faccia diventa obliqua. Con i default del nodo:
\begin{itemize}[leftmargin=1.4em]
  \item esponente di obliquita': \(4\);
  \item cutoff duro: \(65^\circ\);
  \item oltre il cutoff: inflazione a \(10^6\).
\end{itemize}

\begin{figure}[ht]
\centering
\begin{tikzpicture}
\begin{axis}[
  width=0.82\textwidth,
  height=0.34\textwidth,
  xlabel={Obliquita' faccia [gradi]},
  ylabel={Moltiplicatore su $R_{yaw}$},
  xmin=0, xmax=90,
  ymin=1, ymax=1000000,
  ymode=log,
  samples=181,
  grid=both
]
\addplot[thick,blue!70!black,domain=0:65] {1/max((abs(cos(x)))^4,0.0001)};
\addplot[thick,blue!70!black,domain=65:90] {1000000};
\addplot[dashed,black] coordinates {(65,1) (65,1000000)};
\end{axis}
\end{tikzpicture}
\caption{Forma qualitativa-esatta della fiducia nello yaw: fino a \(65^\circ\) il rumore cresce secondo il termine in coseno; oltre \(65^\circ\) il package considera di fatto lo yaw quasi inutilizzabile.}
\end{figure}

\subsubsection{Scalatura della finestra angolare con la distanza}

Nel planner, la finestra angolare effettiva viene scalata con:
\begin{equation}
  \text{dist\_scale} = \min\left(\frac{\text{angular\_window\_ref\_dist}}{\max(\text{range},0.5)}, 2.0\right)
\end{equation}

Con \code{angular\_window\_ref\_dist = 3.0} si ottiene:

\begin{figure}[ht]
\centering
\begin{tikzpicture}
\begin{axis}[
  width=0.82\textwidth,
  height=0.34\textwidth,
  xlabel={Distanza bersaglio [m]},
  ylabel={dist\_scale},
  xmin=0.5, xmax=10,
  ymin=0, ymax=2.1,
  samples=200,
  grid=major
]
\addplot[thick,red!70!black,domain=0.5:10] {min(3/x,2)};
\end{axis}
\end{tikzpicture}
\caption{La finestra angolare e' piu' permissiva a corto raggio, poi si restringe con l'aumentare della distanza.}
\end{figure}

\subsubsection{Clamp del comando assoluto}

Il command controller corrente non filtra il comando con una EMA. Applica invece un clamp sugli angoli assoluti prima di pubblicarli. Per il pitch, con \code{max\_cmd\_angle = 30 deg}, il comportamento ideale e':
\begin{equation}
  \theta_{cmd} = \min(\max(\theta_{plan}, -30^\circ), 30^\circ)
\end{equation}

\begin{figure}[ht]
\centering
\begin{tikzpicture}
\begin{axis}[
  width=0.82\textwidth,
  height=0.34\textwidth,
  xlabel={Pitch richiesto dal piano [gradi]},
  ylabel={Pitch pubblicato [gradi]},
  xmin=-60, xmax=60,
  ymin=-35, ymax=35,
  grid=major
]
\addplot[thick,green!50!black,domain=-60:-30] {-30};
\addplot[thick,green!50!black,domain=-30:30] {x};
\addplot[thick,green!50!black,domain=30:60] {30};
\addplot[dashed,black] coordinates {(-30,-35) (-30,35)};
\addplot[dashed,black] coordinates {(30,-35) (30,35)};
\end{axis}
\end{tikzpicture}
\caption{Il comando segue il piano finche' resta entro i limiti, poi viene saturato. La convergenza dinamica e' responsabilita' del controller gimbal a valle.}
\end{figure}

\subsection{Dettagli di documentazione e coerenza interna del repository}

Nel repository esistono alcuni punti di drift documentale che vale la pena dichiarare esplicitamente nel testo:
\begin{enumerate}[leftmargin=1.4em]
  \item il README del package parla di directory \code{pipeline/}, ma l'albero reale del package usa file radice come \code{auto\_aim\_node.*}, \code{config.hpp}, \code{types.hpp};
  \item il README cita nomi come \code{runtime\_manager}, \code{target\_tracker} e \code{predictor}, mentre il codice consolidato usa \code{tracker\_manager}, \code{tracker} e \code{shot\_planner};
  \item alcuni commenti del launch \code{debug.launch.py} riflettono ancora il vecchio split \code{armor\_tracker + rm\_trajectory}, mentre \pkg\ consolida questi ruoli internamente;
  \item questo non invalida il runtime, ma e' importante per chi legge il repository: documentazione e codice non sono perfettamente allineati nei nomi.
\end{enumerate}

\section{Sezione aggiuntiva: lettura guidata dell'intero codice del package}

\subsection{Come leggere davvero tutto il package}

Leggere ``tutto il codice'' in modo utile non significa copiare ogni riga in ordine alfabetico. Significa ricostruire:
\begin{enumerate}[leftmargin=1.4em]
  \item la funzione di ogni file;
  \item la relazione tra header e implementazione;
  \item il percorso dei dati da ingresso a uscita;
  \item i punti in cui il package conserva stato, decide, scarta o pubblica.
\end{enumerate}

Per questo la lettura seguente e' \emph{file-per-file} e, dentro ogni file, \emph{dichiarazione-per-dichiarazione} o \emph{funzione-per-funzione} quando necessario.

\subsection{Metodo di code review per principianti}

Quando si legge un package ROS~2 grande, il rischio e' perdersi in dettagli locali. Un metodo piu' sicuro e' leggere in tre passaggi.

\paragraph{Primo passaggio: confini.}
Prima di leggere gli algoritmi, individua:
\begin{enumerate}[leftmargin=1.4em]
  \item quali topic entrano;
  \item quali topic escono;
  \item quali parametri cambiano comportamento;
  \item quale callback parte a ogni frame.
\end{enumerate}
Nel package questo significa partire da \code{auto\_aim\_node.hpp/cpp}, non da \code{tracker.cpp}.

\paragraph{Secondo passaggio: tipi dati.}
Poi leggi i tipi di scambio: \code{Observation}, \code{TrackSnapshot}, \code{ShotPlan}, \code{AimDecision}, \code{FireGateResult}, \code{CommandOutput}. Questi tipi sono la grammatica del package. Se sai cosa contiene uno \code{ShotPlan}, leggere \code{ShotPlanner}, \code{EngagementPlanner}, \code{FireGate} e \code{GimbalCommandController} diventa molto piu' semplice.

\paragraph{Terzo passaggio: stato mutabile.}
Infine guarda dove vive lo stato:
\begin{itemize}[leftmargin=1.4em]
  \item nel nodo: mappa dei tracker, id precedenti, stato indirect, timestamp ultimo frame;
  \item nel tracker: EKF, scalar KF, stato della face matched, diagnostica;
  \item nel latency compensator: bias, varianza, warmup;
  \item nel pose adapter: ultima posa, quaternion e freshness.
\end{itemize}

Una review utile non chiede solo ``questa funzione compila?''. Chiede anche: \emph{questa funzione modifica stato? usa timestamp coerenti? scarta dati in modo esplicito? pubblica un output che un altro modulo interpreta con la stessa unita' di misura?}

\subsection{Call graph mentale: da callback a comando}

Il percorso piu' importante da tenere a mente e':
\begin{lstlisting}[style=shellstyle,caption={Call graph semplificato per leggere il runtime}]
armorsCallback
  -> GimbalPoseAdapter::updateForFrame
  -> DetectionConverter::buildObservation
  -> AutoAimNode::processFrame
       -> Trackers::predictAll
       -> Trackers::assignDetections
       -> Trackers::updateMatched
       -> Trackers::spawnFromUnmatched
       -> Trackers::pruneLost
       -> Trackers::collectSnapshots
       -> LatencyCompensator::updateFromFrameStamp
       -> EngagementPlanner::selectBestPlan
            -> ShotPlanner::solveVisibleDirect / solvePredictedDirect / solveIndirect
            -> EngagementPlanner::computePlanScore
       -> FireGate::evaluate
       -> TargetingOutputPublisher::buildAimResult
            -> TargetingDebugPublisher
            -> GimbalCommandController
       -> TargetingOutputPublisher::publish
\end{lstlisting}

Questo call graph e' piu' importante dell'ordine dei file nell'albero. Se sei all'inizio, stampa mentalmente questo percorso e torna a consultarlo ogni volta che una funzione sembra isolata.

\subsection{Checklist di review per ogni blocco}

Questa checklist serve per leggere il codice come un revisore, non solo come un utente.

\paragraph{Measurement.}
Domande da fare mentre leggi \code{detection\_converter.cpp} e \code{pnp\_solver.cpp}:
\begin{enumerate}[leftmargin=1.4em]
  \item cosa succede se \topic{/camera_info} non e' ancora arrivato?
  \item quali detection vengono scartate prima del PnP?
  \item in quale ordine sono i punti 2D e i punti 3D?
  \item quale frame contiene la posa prima e dopo \code{transformToTrackingFrame()}?
  \item una solve OpenCV riuscita viene accettata sempre, o passa anche dal reprojection gate?
\end{enumerate}

\paragraph{Tracking.}
Domande da fare su \code{tracker.cpp} e \code{tracker\_manager.cpp}:
\begin{enumerate}[leftmargin=1.4em]
  \item il tracker sta facendo predict, update o solo coast?
  \item una detection e' associata per id, distanza metrica o Mahalanobis?
  \item quando il codice decide che c'e' un face jump?
  \item quali stati vengono inflazionati, resettati o decoupled?
  \item lo snapshot consegnato al planner contiene dati coerenti con lo stato corrente?
\end{enumerate}

\paragraph{Planning.}
Domande da fare su \code{shot\_planner.cpp} e \code{engagement\_planner.cpp}:
\begin{enumerate}[leftmargin=1.4em]
  \item perche' il planner sceglie visible, predicted o indirect?
  \item il piano e' numericamente valido, meccanicamente reachable e dentro range?
  \item il costo sta penalizzando switch target, switch face, staleness o incertezza?
  \item la scelta del best plan dipende da un singolo target o da un confronto tra tutti gli snapshot?
  \item indirect mode resta stabile vicino alla soglia o oscilla frame per frame?
\end{enumerate}

\paragraph{Fire gate e output.}
Domande da fare su \code{fire\_gate.cpp} e \code{gimbal\_command\_controller.cpp}:
\begin{enumerate}[leftmargin=1.4em]
  \item qual e' il primo blocker che ferma il fuoco?
  \item il blocco dipende da tracking, misura, balistica, range, window o pose source?
  \item il comando finale e' in gradi o radianti?
  \item il comando e' assoluto o relativo?
  \item il consumer downstream legge \code{GimbalCmd} o il \code{Twist} legacy?
\end{enumerate}

Queste domande sono piu' utili di una lettura lineare perche' separano bug di integrazione, bug numerici e scelte intenzionali di design.

\subsection{Mappa completa del package}

L'albero del package comprende:
\begin{itemize}[leftmargin=1.4em]
  \item file radice: \code{README.md}, \code{package.xml}, \code{CMakeLists.txt};
  \item header pubblici radice: \code{auto\_aim\_node.hpp}, \code{config.hpp}, \code{types.hpp}, \code{planning\_types.hpp};
  \item directory \code{measurement/}: \code{armor\_model.hpp}, \code{detection\_converter.hpp/cpp}, \code{pnp\_solver.hpp/cpp};
  \item directory \code{tracking/}: \code{ekf.hpp/cpp}, \code{tracker.hpp/cpp}, \code{tracker\_manager.hpp/cpp}, \code{tracker\_runtime\_factory.hpp/cpp};
  \item directory \code{planning/}: \code{ballistics\_solver.hpp/cpp}, \code{shot\_planner.hpp/cpp}, \code{engagement\_planner.hpp/cpp}, \code{latency\_compensator.hpp/cpp}, \code{fire\_gate.hpp/cpp};
  \item directory \code{io/}: \code{gimbal\_pose\_adapter.hpp/cpp}, \code{gimbal\_command\_controller.hpp/cpp}, \code{debug\_publisher.hpp/cpp}, \code{output\_publisher.hpp/cpp};
  \item entrypoint sorgente: \code{src/auto\_aim\_node.cpp};
  \item test: 8 file \code{gtest} divisi per modulo.
\end{itemize}

\begingroup
\scriptsize
\begin{longtable}{p{0.73\textwidth} p{0.15\textwidth}}
\toprule
\textbf{File} & \textbf{Righe} \\
\midrule
\endfirsthead
\toprule
\textbf{File} & \textbf{Righe} \\
\midrule
\endhead
\code{CMakeLists.txt} & 168 \\
\code{README.md} & 62 \\
\code{package.xml} & 35 \\
\code{include/auto\_aim\_targeting/auto\_aim\_node.hpp} & 111 \\
\code{include/auto\_aim\_targeting/config.hpp} & 116 \\
\code{include/auto\_aim\_targeting/types.hpp} & 106 \\
\code{include/auto\_aim\_targeting/planning\_types.hpp} & 166 \\
\code{include/auto\_aim\_targeting/measurement/armor\_model.hpp} & 70 \\
\code{include/auto\_aim\_targeting/measurement/detection\_converter.hpp} & 58 \\
\code{include/auto\_aim\_targeting/measurement/pnp\_solver.hpp} & 118 \\
\code{include/auto\_aim\_targeting/tracking/ekf.hpp} & 164 \\
\code{include/auto\_aim\_targeting/tracking/tracker.hpp} & 535 \\
\code{include/auto\_aim\_targeting/tracking/tracker\_manager.hpp} & 52 \\
\code{include/auto\_aim\_targeting/tracking/tracker\_runtime\_factory.hpp} & 31 \\
\code{include/auto\_aim\_targeting/planning/ballistics\_solver.hpp} & 46 \\
\code{include/auto\_aim\_targeting/planning/shot\_planner.hpp} & 76 \\
\code{include/auto\_aim\_targeting/planning/engagement\_planner.hpp} & 38 \\
\code{include/auto\_aim\_targeting/planning/latency\_compensator.hpp} & 39 \\
\code{include/auto\_aim\_targeting/planning/fire\_gate.hpp} & 31 \\
\code{include/auto\_aim\_targeting/io/gimbal\_pose\_adapter.hpp} & 59 \\
\code{include/auto\_aim\_targeting/io/gimbal\_command\_controller.hpp} & 39 \\
\code{include/auto\_aim\_targeting/io/debug\_publisher.hpp} & 66 \\
\code{include/auto\_aim\_targeting/io/output\_publisher.hpp} & 70 \\
\code{src/auto\_aim\_node.cpp} & 513 \\
\code{src/measurement/detection\_converter.cpp} & 244 \\
\code{src/measurement/pnp\_solver.cpp} & 235 \\
\code{src/tracking/ekf.cpp} & 559 \\
\code{src/tracking/tracker.cpp} & 1528 \\
\code{src/tracking/tracker\_manager.cpp} & 211 \\
\code{src/tracking/tracker\_runtime\_factory.cpp} & 160 \\
\code{src/planning/ballistics\_solver.cpp} & 106 \\
\code{src/planning/shot\_planner.cpp} & 377 \\
\code{src/planning/engagement\_planner.cpp} & 130 \\
\code{src/planning/latency\_compensator.cpp} & 50 \\
\code{src/planning/fire\_gate.cpp} & 90 \\
\code{src/io/gimbal\_pose\_adapter.cpp} & 194 \\
\code{src/io/gimbal\_command\_controller.cpp} & 69 \\
\code{src/io/debug\_publisher.cpp} & 245 \\
\code{src/io/output\_publisher.cpp} & 162 \\
\code{test/io/test\_gimbal\_command\_controller.cpp} & 70 \\
\code{test/io/test\_gimbal\_pose\_adapter.cpp} & 115 \\
\code{test/pipeline/test\_include\_layout.cpp} & 78 \\
\code{test/planning/test\_engagement\_planner.cpp} & 136 \\
\code{test/planning/test\_fire\_gate.cpp} & 117 \\
\code{test/planning/test\_latency\_compensator.cpp} & 32 \\
\code{test/tracking/test\_runtime\_factory.cpp} & 71 \\
\code{test/tracking/test\_runtime\_manager.cpp} & 163 \\
\bottomrule
\end{longtable}
\endgroup

\subsection{File radice: identita', dipendenze e build}

\subsubsection{\code{README.md}}

Questo file definisce l'identita' del package. Le informazioni principali sono:
\begin{enumerate}[leftmargin=1.4em]
  \item il package e' l'``owner-path targeting runtime'' dello stack di auto-aim;
  \item l'entrypoint operativo e' \code{rm\_auto\_aim::AutoAimNode};
  \item la pipeline viene presentata come: detector -> pose adapter -> detection converter -> trackers -> latency/planner -> fire gate -> output publisher.
\end{enumerate}

Leggendo il README con occhio critico si vedono anche due cose:
\begin{enumerate}[leftmargin=1.4em]
  \item i concetti descritti sono coerenti con il codice reale;
  \item alcuni nomi di file/moduli non sono piu' aggiornati rispetto all'albero reale consolidato.
\end{enumerate}

In pratica il README e' buono per capire il disegno, meno affidabile per cercare il nome esatto di ogni file.

\subsubsection{\code{package.xml}}

Questo file dice che il package:
\begin{itemize}[leftmargin=1.4em]
  \item si chiama \code{auto\_aim\_targeting};
  \item usa \code{ament\_cmake};
  \item dichiara licenza BSD;
  \item dipende da ROS~2 core, TF2, messaggi geometry/sensor/visualization, \code{auto\_aim\_interfaces}, \code{vision\_msgs}, Eigen, angles.
\end{itemize}

Dal punto di vista della lettura del codice, \code{package.xml} serve a capire due cose:
\begin{enumerate}[leftmargin=1.4em]
  \item quali namespace e tipi esterni ci si deve aspettare nei sorgenti;
  \item che il package e' concepito come componente ROS~2, non come semplice libreria standalone.
\end{enumerate}

\subsubsection{\code{CMakeLists.txt}}

Questo file dice esattamente come il package viene assemblato.

\paragraph{Primo blocco.}
Imposta C++17, esporta \code{compile\_commands.json} e abilita \code{-Wall -Werror}. Quindi il progetto e' impostato per essere trattato con warning severi.

Nota pratica per gli editor: \code{CMAKE\_EXPORT\_COMPILE\_COMMANDS} genera il database nel build tree, per esempio \path{build/compile_commands.json}. Se clangd o VS Code mostrano include mancanti mentre \code{colcon build} passa, quasi sempre l'editor non sta usando quel file.

\paragraph{Secondo blocco.}
Chiama \code{find\_package(...)} per tutte le dipendenze. Questo e' importante perche' conferma che:
\begin{itemize}[leftmargin=1.4em]
  \item il PnP usa OpenCV;
  \item il tracking usa Eigen e angles;
  \item l'integrazione spazia su TF2, ROS messages e component framework.
\end{itemize}

\paragraph{Terzo blocco.}
Crea una libreria condivisa con tutti i file \code{src/*.cpp}. L'ordine nell'elenco non cambia il comportamento, ma e' utile per vedere immediatamente quali sono i moduli ``considerati parte del runtime''.

\paragraph{Quarto blocco.}
Registra la libreria come componente ROS~2:
\begin{lstlisting}[style=cppstyle,caption={Dichiarazione del nodo come componente ROS 2 nel build system}]
rclcpp_components_register_node(${PROJECT_NAME}
  PLUGIN rm_auto_aim::AutoAimNode
  EXECUTABLE ${PROJECT_NAME}_node
)
\end{lstlisting}

\paragraph{Quinto blocco.}
Installa target, include directory ed export CMake. Questo rende il package riusabile da altri package del workspace.

\paragraph{Sesto blocco.}
Definisce i test. Questo blocco e' utile perche' mostra la copertura intenzionale:
\begin{itemize}[leftmargin=1.4em]
  \item planner;
  \item latency;
  \item fire gate;
  \item runtime factory;
  \item command controller;
  \item pose adapter;
  \item tracker manager;
  \item include layout.
\end{itemize}

\subsection{Header radice dell'API pubblica}

\subsubsection{\code{auto\_aim\_node.hpp}}

Questo header dichiara la classe \code{AutoAimNode} e quindi l'interfaccia pubblica del nodo.

\paragraph{Struttura logica della classe.}
L'header e' ordinato in modo leggibile:
\begin{enumerate}[leftmargin=1.4em]
  \item runtime path;
  \item leaf callbacks;
  \item inizializzazione e load config;
  \item stato interno.
\end{enumerate}

\paragraph{Cose da notare.}
\begin{enumerate}[leftmargin=1.4em]
  \item \code{processFrame(const Observation \&)} e' il vero cuore del nodo;
  \item esistono callback separati per detection, camera info, micro pose e IMU;
  \item il nodo possiede puntatori unici a tutti i sotto-moduli;
  \item il nodo possiede anche lo stato globale di orchestrazione: \code{dt\_}, \code{trackers\_}, \code{next\_tracker\_id\_}, \code{best\_tracker\_id\_}, \code{previous\_tracker\_id\_}, \code{indirect\_mode\_active\_}.
\end{enumerate}

Questo header da solo chiarisce gia' che il nodo non implementa gli algoritmi nel proprio corpo dati: possiede oggetti specializzati e coordina il loro uso.

\subsubsection{\code{config.hpp}}

Questo file e' la codifica del modello di configurazione del package.

\paragraph{MeasurementConfig.}
Contiene le soglie e i parametri per trasformare bounding box in armature:
\code{max\_armor\_distance}, classi ammesse, \code{light\_ratio}, \code{bbox\_padding\_y}, \code{pnp\_max\_reprojection\_error}, \code{target\_frame}.

\paragraph{TrackingRuntimeConfig.}
Raggruppa:
\begin{itemize}[leftmargin=1.4em]
  \item \code{TrackerConfig};
  \item limiti sul numero di tracker;
  \item rumore di processo globale;
  \item frequenza di refresh.
\end{itemize}

\paragraph{PoseConfig.}
Contiene convenzioni di segno e geometria meccanica del gimbal/camera.

\paragraph{BallisticsConfig.}
Contiene i parametri balistici e i limiti meccanici del pitch.

\paragraph{EngagementConfig.}
Contiene i parametri della logica di ingaggio, della latenza, del fire gate e della cost function. Include anche \code{allow\_fire\_without\_pose}, che rende esplicito quando il sistema puo' sparare senza una pose source reale.

\paragraph{VisualizationConfig.}
Contiene i parametri che in realta' impattano anche i comandi, non solo la visualizzazione: clamp angolari di pitch e yaw.

\paragraph{TargetingConfig.}
E' il contenitore totale. Il metodo \code{ballisticsParams()} converte la config interna nel tipo \code{BallisticsParams} richiesto dal solver balistico.

\subsubsection{\code{types.hpp}}

Questo file definisce i tipi di handoff tra moduli:
\begin{description}[leftmargin=1.8em, style=nextline]
  \item[\code{GimbalPoseState}] yaw, pitch, quaternion, timestamp ultimo update e freshness.
  \item[\code{Observation}] messaggio detector originale, armature convertite, timestamp, stato posa e mappa degli indici originali delle detection.
  \item[\code{FireGateResult}] coppia minimale: \code{fire} + \code{blocker}.
  \item[\code{CommandOutput}] contiene \code{GimbalCmd}, \code{Twist}, \code{Marker}.
  \item[\code{AimResult}] e' il pacchetto finale costruito prima della pubblicazione.
\end{description}

Il metodo inline \code{makePlanningContext()} e' molto importante. Prende configurazione e stato runtime e costruisce il contesto unico letto dal planner. Quindi e' un ponte fondamentale tra tracking e planning.

\subsubsection{\code{planning\_types.hpp}}

Questo file definisce quasi tutto il vocabolario del planner.

\paragraph{Elementi principali.}
\begin{itemize}[leftmargin=1.4em]
  \item \code{AimMode}: enum delle tre modalita' di tiro;
  \item \code{PredictedCenter}: centro del target predetto a tempo futuro;
  \item \code{FaceGeometry}: geometria di una faccia candidata;
  \item \code{BallisticSolution}: output del solver balistico;
  \item \code{ShotPlan}: descrizione completa di un piano di tiro;
  \item \code{CostWeights}: pesi della scoring function;
  \item \code{AimDecision}: coppia \code{TrackSnapshot + ShotPlan};
  \item \code{PlanningContext}: contesto runtime condiviso.
\end{itemize}

\paragraph{Funzioni inline.}
\begin{enumerate}[leftmargin=1.4em]
  \item \code{predictCenter()}: usa posizione, velocita' e accelerazione per predire il centro.
  \item \code{crossRangeSigma()}: proietta la varianza posizionale sulla direzione cross-range rispetto alla bearing line.
\end{enumerate}

Queste due funzioni sono piccole ma concettualmente importanti: la prima e' il ponte con la cinematica; la seconda con la propagazione dell'incertezza verso una quantita' angolare utile al planner.

\subsection{Measurement}

\subsubsection{\code{measurement/armor\_model.hpp}}

Questo file e' piccolo ma fondativo.

\paragraph{Contenuto.}
\begin{enumerate}[leftmargin=1.4em]
  \item enum \code{ArmorType}: SMALL e LARGE;
  \item array \code{ARMOR\_TYPE\_STR};
  \item struct \code{Light};
  \item struct \code{Armor}.
\end{enumerate}

\paragraph{Perche' e' importante.}
Il package separa il concetto di ``armatura interna di pre-PnP'' dal messaggio ROS \code{auto\_aim\_interfaces::msg::Armor}. La struct \code{Armor} di questo file contiene infatti i punti 2D delle light bar e i metadati grezzi del detector. Solo dopo il PnP si ottiene la rappresentazione 3D ROS.

\subsubsection{\code{measurement/pnp\_solver.hpp}}

Questo header e' quasi un mini-tutorial incorporato sul PnP. Dal punto di vista del codice, dichiara:
\begin{itemize}[leftmargin=1.4em]
  \item costruttore con matrice camera e distorsione;
  \item \code{solvePnP()};
  \item \code{solveBestArmorType()};
  \item \code{calculateDistanceToCenter()};
  \item costanti geometriche delle armature;
  \item vettori di punti modello 3D per SMALL e LARGE.
\end{itemize}

\paragraph{Nota di design.}
L'header esplicita gia' che si usa IPPE per un oggetto planare. Questo e' un raro caso in cui i commenti rendono il codice significativamente piu' leggibile.

\subsubsection{\code{measurement/pnp\_solver.cpp}}

Questo file implementa tre responsabilita':
\begin{enumerate}[leftmargin=1.4em]
  \item costruzione del modello 3D della piastra;
  \item risoluzione PnP;
  \item misura della distanza dal centro immagine.
\end{enumerate}

\paragraph{Costruttore.}
Fa copie profonde di matrice camera e distorsione, poi pre-calcola i quattro punti 3D per armatura piccola e grande. Quindi il costo di costruzione del modello non viene ripagato a ogni frame.

\paragraph{\code{solvePnP()}.}
Passo per passo:
\begin{enumerate}[leftmargin=1.4em]
  \item raccoglie i 4 punti 2D in ordine coerente con i 4 punti 3D;
  \item seleziona il modello 3D piccolo o grande;
  \item chiama \code{cv::solvePnP(..., SOLVEPNP\_IPPE)}.
\end{enumerate}

\paragraph{\code{solveBestArmorType()}.}
Prova sia il modello SMALL sia il modello LARGE. Per ogni solve riuscita calcola l'errore di reproiezione medio con \code{computeReprojectionError()} e sceglie il tipo con errore minore. Questo e' un miglioramento importante per robustezza: la scelta del modello fisico avviene confrontando quanto bene il modello spiega i punti immagine.

\paragraph{\code{computeReprojectionError()}.}
Proietta i punti modello con la posa stimata e confronta i punti proiettati con i punti 2D originali. Restituisce infinito se incontra numeri non finiti o array incoerenti. Questa funzione rende possibile un gate quantitativo sulla qualita' della soluzione PnP.

\paragraph{\code{calculateDistanceToCenter()}.}
Legge il principal point \((c_x, c_y)\) dalla matrice camera e restituisce la norma euclidea in pixel rispetto al centro immagine. Questa quantita' viene usata per privilegiare detection vicine al centro in fase di init del tracker.

\subsubsection{\code{measurement/detection\_converter.hpp}}

L'header dichiara l'oggetto che traduce detector output in \code{Observation}. Le funzioni private mostrano gia' l'intera pipeline:
\begin{enumerate}[leftmargin=1.4em]
  \item \code{detectArmors()};
  \item \code{buildArmorFromDetection()};
  \item \code{solveArmorPose()};
  \item \code{transformToTrackingFrame()};
  \item \code{filterImplausibleDetections()}.
\end{enumerate}

\subsubsection{\code{measurement/detection\_converter.cpp}}

Questo e' il file operativo del blocco di misura.

\paragraph{Costruttore.}
Salva config, TF buffer e logger. Non crea subito il PnP solver.

\paragraph{\code{setCameraInfo()}.}
Crea il \code{PnPSolver}. Questo e' il punto in cui la measurement pipeline diventa realmente pronta.

\paragraph{\code{buildObservation()}.}
E' la funzione centrale. Crea un \code{Observation}, copia il messaggio originale, salva timestamp e pose state, poi:
\begin{enumerate}[leftmargin=1.4em]
  \item estrae armature dal detector;
  \item trasforma le pose al frame di tracking;
  \item filtra le detection implausibili.
\end{enumerate}

\paragraph{\code{detectArmors()}.}
Itera su tutte le detection e per ciascuna prova:
\begin{enumerate}[leftmargin=1.4em]
  \item \code{buildArmorFromDetection()};
  \item \code{solveArmorPose()}.
\end{enumerate}
Solo le detection che passano entrambi i passi entrano nell'array finale.

\paragraph{\code{buildArmorFromDetection()}.}
Esegue la costruzione della struct \code{Armor} interna:
\begin{enumerate}[leftmargin=1.4em]
  \item legge bbox;
  \item corregge il centro \(y\) sottraendo \code{bbox\_padding\_y};
  \item scarta larghezza/altezza \(<1\);
  \item scarta \code{light\_ratio <= 0};
  \item filtra per class id;
  \item calcola punti top/bottom delle light bar usando \code{light\_ratio}.
\end{enumerate}

\paragraph{\code{solveArmorPose()}.}
Chiama \code{solveBestArmorType()}, scarta soluzioni con reprojection error superiore a \code{pnp.max\_reprojection\_error}, riempie il messaggio ROS \code{Armor}, converte la rotazione con Rodrigues e scarta quaternion degeneri.

\paragraph{\code{transformToTrackingFrame()}.}
Se il lookup TF fallisce, svuota completamente le armature del frame. Questa e' una scelta drastica ma sicura: evita di mischiare pose in frame sbagliati con pose corrette.

\paragraph{\code{filterImplausibleDetections()}.}
Scorre l'array e cancella in place pose con quota assurda o troppo lontane. Cancella in parallelo anche gli indici delle detection originali.

\subsection{Tracking}

\subsubsection{\code{tracking/ekf.hpp}}

Questo header definisce una classe EKF generica tramite dependency injection. Le idee chiave dichiarate qui sono:
\begin{enumerate}[leftmargin=1.4em]
  \item il filtro non conosce nulla di armature o RoboMaster;
  \item il comportamento dipende da funzioni esterne \(f, h, J_f, J_h, Q, R\);
  \item l'EKF espone utility importanti oltre a predict/update: \code{mahalanobis}, \code{mahalanobisAt}, \code{inflateCovariance}, \code{decoupleState}, \code{syncPrior}, \code{updateWithModel}.
\end{enumerate}

Dal punto di vista di design, questa e' una scelta molto pulita: il filtro e' riusabile e il dominio resta nella factory del tracker.

\subsubsection{\code{tracking/ekf.cpp}}

Questo file implementa la matematica dell'EKF.

\paragraph{Costruttore.}
Riceve tutte le funzioni e inizializza \(P_{pri}\), \(P_{post}\), \(P_0\), identita' e vettori di stato.

\paragraph{\code{setState()}.}
Scrive contemporaneamente \code{x\_post} e \code{x\_pri}. Questo e' importante per evitare inconsistenze quando il tracker forza manualmente uno snap o un reset.

\paragraph{\code{resetCovariance()}.}
Ripristina \(P\) alla covarianza iniziale. Serve quando la stima e' talmente incoerente da richiedere una ripartenza probabilistica larga.

\paragraph{\code{syncPrior()}.}
Copia posterior in prior. Viene usato dopo modifiche manuali allo stato, per fare in modo che il prossimo update si linearizzi attorno allo stato corretto e non a quello precedente.

\paragraph{\code{inflateCovariance()}.}
Scala riga e colonna di uno stato per inflazionarne l'incertezza mantenendo la struttura correlazionale il piu' possibile coerente.

\paragraph{\code{decoupleState()}.}
E' fondamentale per il raggio. Azzera le cross-covarianze dello stato selezionato e ne imposta la varianza desiderata. Nel package viene usato per impedire all'EKF 9D di ``combattere'' contro lo scalar KF del raggio.

\paragraph{\code{predict()}.}
Esegue:
\begin{enumerate}[leftmargin=1.4em]
  \item linearizzazione con \(F\);
  \item aggiornamento \(Q\);
  \item prediction di stato;
  \item prediction di covarianza;
  \item clamp opzionale della diagonale di \(P\);
  \item controllo spettrale sul condition number;
  \item copia prior -> posterior per consentire coast.
\end{enumerate}

\paragraph{\code{mahalanobis()} e \code{mahalanobisAt()}.}
Calcolano la distanza statistica della misura dal prior. La versione \code{At()} permette di linearizzare \(H\) in un punto diverso dal prior, utile nei face jump.

\paragraph{\code{update()}.}
Esegue il classico update EKF con Joseph form e fallback sicuro a prediction-only se \(S\) non e' positiva definita.

\paragraph{\code{updateWithModel()}.}
Permette un update sequenziale con modello osservativo esterno. Nel package serve per la secondary face fusion.

\subsubsection{\code{tracking/tracker.hpp}}

Questo header e' il piu' denso del package.

\paragraph{Tipi dichiarati.}
\begin{description}[leftmargin=1.8em, style=nextline]
  \item[\code{FaceBinding}] rappresenta l'identita' di una faccia associata a un tracker.
  \item[\code{TrackSnapshot}] e' la fotografia immutable del tracker data al planner.
  \item[\code{ScalarKF}] filtro 1D per i raggi.
  \item[\code{ArmorsNum}] enum per 2, 3, 4 armature.
  \item[\code{TrackerConfig}] enorme struttura con tutti i parametri di tracking.
\end{description}

\paragraph{Classe \code{Tracker}.}
L'header espone:
\begin{itemize}[leftmargin=1.4em]
  \item metodi di init;
  \item ciclo predict/update;
  \item utility per snapshot e prediction di facce;
  \item molti campi pubblici usati come stato diagnostico o come bridge con i publisher;
  \item molte utility private per jump, face resolution, stationary handling.
\end{itemize}

Il fatto che molti campi siano pubblici puo' sembrare poco ``OO'', ma qui e' una scelta pragmatica: il tracker e' un oggetto di runtime molto interno, e il package privilegia l'accessibilita' del suo stato diagnostico.

\subsubsection{\code{tracking/tracker.cpp}}

Questo e' il file piu' importante del package insieme al nodo principale.

\paragraph{Costruttore.}
Mette il tracker in stato LOST, inizializza vettori e imposta \code{max\_match\_yaw\_diff\_} a \(60^\circ\).

\paragraph{\code{init()}.}
Prende un array di armature, sceglie quella piu' vicina al centro immagine, inizializza EKF e stato del tracker. L'idea e' minimizzare l'errore di bootstrap scegliendo la detection piu' affidabile.

\paragraph{\code{initFromArmor()}.}
Versione diretta di bootstrap da una sola \code{Armor}. E' quella usata in \code{spawnFromUnmatched()}.

\paragraph{\code{update()}.}
E' la funzione gigantesca che realizza quasi tutto il comportamento del tracker. Le sue fasi sono:
\begin{enumerate}[leftmargin=1.4em]
  \item predict interna o riuso della prediction gia' fatta;
  \item estrazione delle armature con lo stesso \code{tracked\_id};
  \item selezione della detection primaria e, se c'e', di un'altra faccia utile;
  \item calcolo di errori di posizione, yaw e obliquita';
  \item classificazione match / jump / miss;
  \item update EKF in caso di match normale;
  \item adattamento di raggio e \code{dz};
  \item eventuale secondary face fusion;
  \item handleArmorJump in caso di jump;
  \item clamp di sicurezza su raggio e \code{v\_yaw};
  \item avanzamento della state machine;
  \item drop del tracker se fuori range;
  \item normalizzazione finale dello yaw.
\end{enumerate}

\paragraph{Perche' \code{update()} e' lunga.}
Non per mancanza di struttura, ma perche' contiene la semantica di runtime di un target: stima, diagnostica, classificazione, gestione geometria e state machine.

\paragraph{\code{computeAdaptiveDamping()}.}
Modula il damping di velocita' in base a:
\begin{itemize}[leftmargin=1.4em]
  \item overshoot lineare;
  \item overshoot angolare;
  \item stato TEMP\_LOST;
  \item stationary detection.
\end{itemize}

In pratica il tracker cambia quanto si fida della propria velocita' a seconda del comportamento osservato.

\paragraph{\code{predictStep()}.}
Esegue predict dell'EKF e dei due scalar KF, sincronizza il raggio nel vettore di stato e memorizza le velocita' prior per il calcolo successivo dell'accelerazione.

\paragraph{\code{computeAssignmentCost()}.}
Costruisce una misura \([x,y,z,yaw]\) senza alterare lo stato di unwrapping e ritorna la Mahalanobis. Serve solo all'assegnazione globale.

\paragraph{\code{initEKF()}.}
Back-calcola il centro dal primo armor pose e inizializza anche le due stime di raggio e \code{dz}.

\paragraph{\code{updateArmorsNum()}.}
Attualmente forza sempre \code{NORMAL\_4}. E' importante dirlo perche' il codice ha la struttura per piu' casi, ma la logica runtime non li usa ancora.

\paragraph{\code{handleArmorJump()}.}
Gestisce il cambio faccia. Esegue:
\begin{enumerate}[leftmargin=1.4em]
  \item detection di inversione del verso di spin;
  \item snap dello yaw;
  \item eventuale pair switch con swap dei raggi e aggiornamento \code{dz};
  \item detection di divergenza del centro e hard reset se necessario;
  \item inflazione della covarianza degli stati ``snappati'';
  \item syncPrior;
  \item decouple del raggio;
  \item update EKF con la misura corrente;
  \item aggiornamento del raggio attivo usando la nuova faccia.
\end{enumerate}

\paragraph{\code{orientationToYaw()}.}
Converte quaternion -> yaw continuo. E' una funzione piccola ma critica per evitare salti numerici di \(2\pi\) nella velocita' angolare.

\paragraph{\code{getArmorPositionFromState()}.}
Implementa esattamente l'inverso del modello osservativo: da centro + yaw + raggio a posa della faccia.

\paragraph{\code{updateSmoothedAcceleration()}.}
Calcola un'accelerazione ``innovation-based'' a partire dalla differenza tra velocita' posterior e prior, correggendo anche l'effetto del damping. Poi la filtra con una EMA.

\paragraph{\code{snapshot()}.}
Serializza tutto lo stato utile del tracker in un \code{TrackSnapshot}. Questa funzione e' il vero confine tra tracking e planning.

\paragraph{\code{predictCenter()}, \code{predictFace()}, \code{predictMatchedFace()}.}
Sono utility geometriche per il planner:
\begin{itemize}[leftmargin=1.4em]
  \item predicono centro futuro;
  \item predicono facce candidate;
  \item calcolano visibilita' e osservabilita'.
\end{itemize}

\paragraph{\code{buildFaceBindingFromState()}.}
Costruisce la rappresentazione geometrica di una faccia a partire dallo stato corrente e dal suo indice.

\paragraph{\code{updateFaceContext()}.}
Risoluzione della faccia misurata rispetto allo stato. Di fatto attribuisce un'identita' alla faccia matched.

\paragraph{\code{updateStationaryState()} e \code{resetStationaryState()}.}
Gestiscono la logica di bersaglio fermo e la sua confidenza.

\paragraph{\code{resolveMatchedFace()}.}
Utility globale, non metodo membro. Scorre tutte le facce teoriche e trova quella piu' coerente con misura e stato, popolando un \code{FaceBinding}.

\subsubsection{\code{tracking/tracker\_manager.hpp}}

L'header dichiara la classe \code{Trackers} come contenitore di funzioni statiche sullo store di tracker. Non c'e' uno stato interno nella classe; lo stato e' tutto nella mappa esterna.

\subsubsection{\code{tracking/tracker\_manager.cpp}}

Questo file e' il coordinatore multi-target.

\paragraph{\code{predictAll()}.}
Aggiorna \code{lost\_thres}, damping e prediction per tutti i tracker attivi.

\paragraph{\code{assignDetections()}.}
Per ogni detection cerca il tracker compatibile con stesso \code{tracked\_id} e Mahalanobis minima.

\paragraph{\code{updateMatched()}.}
Costruisce per ogni tracker il suo sotto-array di detection assegnate e chiama \code{tracker.update()}.

\paragraph{\code{spawnFromUnmatched()}.}
Per ogni detection non assegnata:
\begin{enumerate}[leftmargin=1.4em]
  \item verifica che non sia troppo vicina a tracker gia' esistenti;
  \item verifica il cap sul numero di tracker;
  \item crea un nuovo tracker via factory;
  \item assegna \code{tracker\_id};
  \item salva timestamp ultima misura.
\end{enumerate}

\paragraph{\code{pruneLost()}.}
Rimuove dalla mappa i tracker in stato LOST.

\paragraph{\code{findTracker()} e \code{collectSnapshots()}.}
Sono utility semplici: la seconda aggiorna anche l'accelerazione smussata prima di serializzare.

\subsubsection{\code{tracking/tracker\_runtime\_factory.hpp}}

Header minimalista. Dichiara una factory che possiede config e \code{DtProvider}. Il provider di \(dt\) viene passato come lambda, cosi' ogni tracker puo' leggere il \code{dt\_} attuale del nodo senza doverne duplicare la gestione.

\subsubsection{\code{tracking/tracker\_runtime\_factory.cpp}}

Questo file costruisce un tracker gia' completamente cablato.

\paragraph{Cosa crea.}
\begin{enumerate}[leftmargin=1.4em]
  \item il tracker;
  \item la lambda di processo \(f\);
  \item il Jacobiano \(J_f\);
  \item il modello osservativo \(h\);
  \item il Jacobiano \(J_h\);
  \item il generatore di \(Q\);
  \item il generatore di \(R\);
  \item la covarianza iniziale \(P_0\);
  \item i limiti di covarianza e il max condition number.
\end{enumerate}

\paragraph{Idea chiave.}
Il filtro non e' configurato dentro il costruttore del tracker, ma dalla factory. Quindi la classe \code{Tracker} resta il ``motore logico'' e la factory il luogo in cui viene deciso il modello matematico.

\subsection{Planning}

\subsubsection{\code{planning/ballistics\_solver.hpp}}

Header molto piccolo. Definisce:
\begin{itemize}[leftmargin=1.4em]
  \item \code{BallisticsParams};
  \item classe \code{BallisticsSolver} con il solo metodo \code{solve()}.
\end{itemize}

\subsubsection{\code{planning/ballistics\_solver.cpp}}

Questo file implementa il solver iterativo della traiettoria.

\paragraph{Passi di \code{solve()}.}
\begin{enumerate}[leftmargin=1.4em]
  \item scarta casi degeneri;
  \item usa \(t = d/v\) come seed iniziale;
  \item per un massimo di 12 iterazioni:
  \begin{itemize}[leftmargin=1.4em]
    \item stima la caduta gravitazionale;
    \item ricava il pitch richiesto;
    \item clampa il pitch ai limiti meccanici;
    \item ricalcola il tempo di volo orizzontale con drag lineare o quadratico;
    \item interrompe se converge.
  \end{itemize}
  \item infine marca \code{reachable} e \code{valid}.
\end{enumerate}

\subsubsection{\code{planning/shot\_planner.hpp}}

L'header dichiara l'oggetto che costruisce i singoli \code{ShotPlan}:
\begin{itemize}[leftmargin=1.4em]
  \item \code{solveVisibleDirect()};
  \item \code{solvePredictedDirect()};
  \item \code{solveIndirect()}.
\end{itemize}

Dichiara anche helper cruciali: \code{buildFaceGeometry()}, \code{isBetterCandidate()}, \code{nextAlignmentTime()}, \code{buildPlan()}.

\subsubsection{\code{planning/shot\_planner.cpp}}

Questo file implementa tutta la logica di costruzione dei candidati di tiro.

\paragraph{\code{solveVisibleDirect()}.}
Usa la faccia matched del tracker, itera fino a 3 volte per rendere coerenti prediction time e flight time, e costruisce un piano se la balistica resta valida.

\paragraph{\code{solvePredictedDirect()}.}
Esamina tutte le facce teoriche del target, costruisce un candidato per ciascuna e tiene la migliore secondo \code{isBetterCandidate()}.

\paragraph{\code{solveIndirect()}.}
Per target che ruotano molto:
\begin{enumerate}[leftmargin=1.4em]
  \item calcola la periodicita' della rotazione;
  \item stima i prossimi istanti di allineamento per ogni faccia;
  \item genera candidati su piu' occorrenze future;
  \item confronta tempo di arrivo del proiettile e tempo di allineamento;
  \item costruisce un margine di firing utile all'indirect mode.
\end{enumerate}

\paragraph{\code{isBetterCandidate()}.}
E' un comparatore multi-criterio. Per direct mode il criterio primario e' \code{fire\_window\_margin}; per indirect mode e' il \code{timing\_residual}. Poi usa reachability, criterio secondario e range come tie-breaker.

\paragraph{\code{buildFaceGeometry()}.}
Converte snapshot + faccia + tempo predetto in una geometria completa della faccia: posizione, quota, range, bearing, visibilita'.

\paragraph{\code{nextAlignmentTime()}.}
Dato yaw del corpo, velocita' di yaw e face spacing, restituisce il prossimo istante in cui una faccia si allineera' alla bearing desiderata.

\paragraph{\code{buildPlan()}.}
E' il costruttore finale del piano. Qui il package combina:
\begin{itemize}[leftmargin=1.4em]
  \item geometria della faccia;
  \item solver balistico;
  \item relative yaw/pitch rispetto al gimbal;
  \item stima dell'incertezza angolare;
  \item margine di fuoco.
\end{itemize}

\subsubsection{\code{planning/engagement\_planner.hpp}}

Header molto semplice: dichiara l'oggetto che seleziona il miglior piano tra tutti i target.

\subsubsection{\code{planning/engagement\_planner.cpp}}

\paragraph{\code{selectBestPlan()}.}
Per ogni snapshot:
\begin{enumerate}[leftmargin=1.4em]
  \item scarta target non in tracking;
  \item decide se usare indirect mode in base a \(|v\_yaw|\), soglia e storico del target precedente;
  \item prova prima visible direct se possibile;
  \item altrimenti usa predicted direct o indirect;
  \item calcola measurement age effettiva;
  \item scarta piani senza balistica valida o fuori range;
  \item assegna score;
  \item conserva il best score minimo.
\end{enumerate}

Il dettaglio di isteresi e' importante: se il target corrente e' lo stesso del frame precedente e l'indirect mode era gia' attiva, il planner resta in indirect finche' \(|v\_yaw|\) non scende sotto il \(70\%\) della soglia. Questo evita l'alternanza rapida direct/indirect quando la velocita' angolare oscilla vicino alla soglia.

\paragraph{\code{computePlanScore()}.}
Costruisce il costo totale a partire da:
\begin{itemize}[leftmargin=1.4em]
  \item slew normalizzato rispetto al budget di velocita' del gimbal;
  \item incertezza;
  \item visibilita';
  \item staleness;
  \item penalita' extra per temp lost, switch target, switch face, margine negativo e piano irraggiungibile.
\end{itemize}

\subsubsection{\code{planning/latency\_compensator.hpp}}

Header essenziale: stato interno minimo, costruttore e tre metodi. Dice gia' chiaramente che la latenza e' modellata come stima adattiva, non come costante fissa.

\subsubsection{\code{planning/latency\_compensator.cpp}}

\paragraph{\code{updateFromFrameStamp()}.}
Misura latenza come \((now - frame\_stamp)\), scarta valori implausibili, applica warmup e gating statistico, poi aggiorna bias e varianza con EMA.

\paragraph{\code{transportDelay()}.}
Restituisce semplicemente il ritardo grezzo non negativo del frame.

\subsubsection{\code{planning/fire\_gate.hpp}}

Header minimale: classe con config e metodo \code{evaluate()}.

\subsubsection{\code{planning/fire\_gate.cpp}}

Questo file e' quasi autoesplicativo. Fa una catena di controlli ordinati e ritorna al primo blocker:
\begin{enumerate}[leftmargin=1.4em]
  \item not tracking;
  \item temp lost;
  \item stale measurement;
  \item invalid ballistics;
  \item unreachable pitch;
  \item out of range;
  \item negative window;
  \item stale pose se una pose source e' configurata;
  \item no pose source se \code{pose\_source=none} e \code{allow\_fire\_without\_pose=false}.
\end{enumerate}

Il codice non fa piu' un gate diretto su \code{relative\_yaw} o \code{relative\_pitch}: il commento del file chiarisce che la convergenza del gimbal viene demandata al controller downstream. Il gate principale resta \code{fire\_window\_margin}, che incorpora finestra angolare, distanza, incertezza e obliquita'.

\subsection{IO}

\subsubsection{\code{io/gimbal\_pose\_adapter.hpp}}

Header che dichiara l'adattatore tra sorgenti di posa e stato consumabile dal resto del package. Espone callback separati per micro pose e IMU, piu' metodi di query e TF broadcast.

\subsubsection{\code{io/gimbal\_pose\_adapter.cpp}}

\paragraph{\code{onMicroPose()}.}
Legge pitch e yaw dai campi posizione del messaggio, applica segni configurabili, costruisce un quaternion RPY e aggiorna il tempo dell'ultima posa valida.

\paragraph{\code{onCameraImu()}.}
Legge direttamente il quaternion dell'IMU, ne estrae yaw/pitch/roll e salva tutto nello stato interno.

\paragraph{\code{updateForFrame()}.}
Controlla la freschezza rispetto allo stamp del frame e, se necessario, emette warning throttled. In ogni caso ribroadcasta il TF camera.

\paragraph{\code{currentPoseState()} e \code{isFresh()}.}
Restituiscono la fotografia dello stato interno. \code{isFresh()} ha la logica particolare per \code{pose\_source=none}.

\paragraph{\code{broadcastCameraTF()}.}
Combina:
\begin{itemize}[leftmargin=1.4em]
  \item quaternion del gimbal;
  \item quaternion di convenzione ottica;
  \item pivot height;
  \item offset camera ruotato.
\end{itemize}
Poi inserisce il TF nel buffer e lo trasmette.

\subsubsection{\code{io/gimbal\_command\_controller.hpp}}

Header piccolo. Dichiara l'oggetto che si occupa di hold output, shot output e clamp angolari. Non contiene piu' stato di smoothing.

\subsubsection{\code{io/gimbal\_command\_controller.cpp}}

\paragraph{\code{makeHoldOutput()}.}
Pubblica comando nullo e marker di delete.

\paragraph{\code{makeShotOutput()}.}
Passi:
\begin{enumerate}[leftmargin=1.4em]
  \item clamp di yaw/pitch assoluti;
  \item conversione in gradi per \code{GimbalCmd};
  \item costruzione del \code{Twist} di compatibilita';
  \item costruzione del marker sferico del punto d'impatto.
\end{enumerate}

Il \code{Twist} usa radianti nei campi \code{angular.y} e \code{angular.z}; il messaggio principale \code{GimbalCmd} usa gradi e contiene anche \code{fire\_cmd}. Per integrazione esterna conviene usare \code{GimbalCmd} come fonte primaria.

\subsubsection{\code{io/debug\_publisher.hpp}}

Header che dichiara l'oggetto di formattazione per:
\begin{itemize}[leftmargin=1.4em]
  \item \code{Target};
  \item \code{TrackerInfo};
  \item bounding box ottimale;
  \item marker array.
\end{itemize}

\subsubsection{\code{io/debug\_publisher.cpp}}

\paragraph{\code{fillTargetMsg()}.}
Riempie il messaggio \code{Target} da un tracker. Gestisce esplicitamente il caso DETECTING e i casi non TRACKING/TEMP\_LOST pubblicando un target vuoto.

\paragraph{\code{makeTrackerInfo()}.}
Traduce errori diagnostici del tracker in un messaggio piccolo e leggibile.

\paragraph{\code{selectOptimalBBox()}.}
Trova la detection 2D originale piu' vicina alla \code{tracked\_armor} del best tracker e la restituisce come bbox ottimale.

\paragraph{\code{buildMarkers()}.}
Costruisce o marker di tracking o marker di delete, poi aggiunge sempre i marker principali di posizione e velocita'.

\paragraph{\code{buildTrackingMarkers()}.}
Costruisce:
\begin{itemize}[leftmargin=1.4em]
  \item marker del centro;
  \item freccia di velocita' lineare;
  \item freccia di velocita' angolare;
  \item marker per ciascuna armatura del target.
\end{itemize}

\paragraph{\code{buildDeleteMarkers()}.}
Pubblica marker di delete per ripulire RViz quando il target non e' piu' attivo.

\subsubsection{\code{io/output\_publisher.hpp}}

Header che dichiara l'aggregatore finale dell'output. Possiede sia il \code{TargetingDebugPublisher} sia il \code{GimbalCommandController}, oltre ai publisher ROS.

\subsubsection{\code{io/output\_publisher.cpp}}

\paragraph{Costruttore.}
Inizializza i template dei marker e costruisce i sotto-oggetti interni.

\paragraph{\code{buildAimResult()}.}
E' la funzione che:
\begin{enumerate}[leftmargin=1.4em]
  \item serializza tutti i tracker in \code{Targets};
  \item costruisce il best \code{Target};
  \item prepara \code{TrackerInfo};
  \item seleziona la bbox ottimale;
  \item costruisce marker debug;
  \item decide se generare hold output o shot output.
\end{enumerate}

\paragraph{\code{publish()}.}
Pubblica tutti i messaggi sulle rispettive interfacce ROS.

\subsection{Nodo principale}

\subsubsection{\code{src/auto\_aim\_node.cpp}}

Questo file collega tutti i moduli letti fin qui.

\paragraph{Costruttore.}
Chiama in sequenza:
\begin{enumerate}[leftmargin=1.4em]
  \item \code{loadConfig()};
  \item \code{createRosInterfaces()};
  \item \code{buildPipeline()}.
\end{enumerate}

\paragraph{\code{armorsCallback()}.}
Verifica che la pipeline sia pronta, aggiorna lo stato posa, costruisce un \code{Observation} e lo passa a \code{processFrame()}.

\paragraph{\code{processFrame()}.}
E' la regia completa del runtime del package. E' il posto in cui measurement, tracking, planning, fire gate e publishing vengono concatenati. Oltre alla sequenza algoritmica, fa anche tre cose difensive:
\begin{enumerate}[leftmargin=1.4em]
  \item prende un mutex sullo stato tracker;
  \item registra che almeno un frame e' arrivato, per il watchdog;
  \item costruisce messaggi vuoti sicuri se \code{armors} o \code{detection\_msg} sono nulli.
\end{enumerate}

\paragraph{\code{detectorWatchdogCallback()} e \code{publishHoldOutput()}.}
Sono il percorso di sicurezza quando il detector smette di pubblicare. Se il timeout scade, il nodo non prova a pianificare su dati vecchi: pubblica un hold output, resetta best tracker e resetta anche \code{previous\_tracker\_id}, \code{previous\_face\_index} e \code{indirect\_mode\_active}.

\paragraph{\code{updateTimestep()}.}
Aggiorna \code{dt\_} clampando l'intervallo. Piccola funzione, grande impatto numerico.

\paragraph{\code{isPipelineReady()}.}
Controlla che tutti i sotto-oggetti essenziali siano stati creati.

\paragraph{Leaf callbacks.}
\begin{itemize}[leftmargin=1.4em]
  \item \code{cameraInfoCallback()}: inizializza PnP;
  \item \code{microPoseCallback()}: inoltra al pose adapter;
  \item \code{cameraImuCallback()}: inoltra al pose adapter.
\end{itemize}

\paragraph{Load config.}
Le funzioni \code{loadMeasurementConfig()}, \code{loadTrackingConfig()}, \code{loadPoseConfig()}, \code{loadBallisticsConfig()}, \code{loadEngagementConfig()}, \code{loadVisualizationConfig()} dichiarano e leggono tutti i parametri ROS.

Due dettagli meritano attenzione: \code{loadMeasurementConfig()} rimuove esplicitamente la classe \code{"1"} dall'insieme \code{target\_classes}; \code{loadEngagementConfig()} dichiara i pesi \code{cost.*}, quindi il launch puo' realmente modificare la funzione di costo.

\paragraph{Create interfaces.}
\begin{itemize}[leftmargin=1.4em]
  \item \code{createServices()}: crea \topic{/tracker/reset};
  \item \code{createTfInfrastructure()}: prepara buffer/listener/broadcaster;
  \item \code{createPublishers()}: crea tutti i publisher e costruisce \code{TargetingOutputPublisher};
  \item \code{createSubscriptions()}: sottoscrive camera info, detector e opzionalmente micro pose o IMU;
  \item timer watchdog: ogni 50 ms controlla se il detector e' fermo.
\end{itemize}

\paragraph{\code{buildPipeline()}.}
Istanzia i moduli veri e propri nell'ordine logico:
\begin{enumerate}[leftmargin=1.4em]
  \item \code{DetectionConverter};
  \item \code{GimbalPoseAdapter};
  \item \code{TrackerFactory};
  \item \code{LatencyCompensator};
  \item \code{FireGate};
  \item \code{BallisticsSolver};
  \item \code{ShotPlanner};
  \item \code{EngagementPlanner}.
\end{enumerate}

\subsection{Test}

\subsubsection{\code{test/io/test\_gimbal\_command\_controller.cpp}}

Contiene tre test:
\begin{enumerate}[leftmargin=1.4em]
  \item verifica che l'hold output resetti comando e marker;
  \item verifica che lo shot output usi angoli assoluti;
  \item verifica il clamp di yaw e pitch.
\end{enumerate}

\subsubsection{\code{test/io/test\_gimbal\_pose\_adapter.cpp}}

Contiene una fixture ROS~2 e due test:
\begin{enumerate}[leftmargin=1.4em]
  \item verifica che a posa neutra il TF includa altezza pivot e offset camera;
  \item verifica che ruotando lo yaw ruoti anche l'offset camera nel TF risultante.
\end{enumerate}

\subsubsection{\code{test/pipeline/test\_include\_layout.cpp}}

Verifica:
\begin{itemize}[leftmargin=1.4em]
  \item che la directory \code{include/auto\_aim\_targeting} contenga esattamente le directory e i file attesi;
  \item che gli include flat siano ammessi solo per alcuni header radice.
\end{itemize}

E' un test di igiene dell'API pubblica, non di algoritmo.

\subsubsection{\code{test/planning/test\_engagement\_planner.cpp}}

Costruisce snapshot sintetici e verifica:
\begin{enumerate}[leftmargin=1.4em]
  \item preferenza per \code{VISIBLE\_DIRECT} se la faccia matched e' fresca;
  \item fallback a \code{PREDICTED\_DIRECT} se la faccia e' stale;
  \item uso di \code{INDIRECT} per target con spin elevato.
\end{enumerate}

\subsubsection{\code{test/planning/test\_fire\_gate.cpp}}

Verifica:
\begin{itemize}[leftmargin=1.4em]
  \item blocco per stale measurement;
  \item permesso di fuoco quando tutte le condizioni sono soddisfatte;
  \item blocco quando manca una pose source e \code{allow\_fire\_without\_pose=false};
  \item permesso esplicito senza pose source quando il parametro di debug lo consente;
  \item blocco per margine negativo, TEMP\_LOST e posa stale quando una pose source esiste.
\end{itemize}

\subsubsection{\code{test/planning/test\_latency\_compensator.cpp}}

Verifica due aspetti:
\begin{enumerate}[leftmargin=1.4em]
  \item che il bias venga aggiornato durante warmup;
  \item che il transport delay non diventi negativo.
\end{enumerate}

\subsubsection{\code{test/tracking/test\_runtime\_factory.cpp}}

Verifica che la factory costruisca un tracker con parametri EKF coerenti con la configurazione passata.

\subsubsection{\code{test/tracking/test\_runtime\_manager.cpp}}

Verifica tre aspetti:
\begin{enumerate}[leftmargin=1.4em]
  \item \code{resolveMatchedFace()} mantiene l'identita' corretta della faccia misurata;
  \item \code{predictMatchedFace()} conserva l'identita' della faccia storica;
  \item \code{collectSnapshots()} ignora i tracker LOST.
\end{enumerate}

\subsection{Cosa insegna la lettura completa del codice}

Dopo aver letto file per file tutto il package, emergono alcune conclusioni forti:
\begin{enumerate}[leftmargin=1.4em]
  \item il package e' scritto come pipeline modulare, ma con un forte stato condiviso orchestrato dal nodo principale;
  \item il tracking e' la parte piu' sofisticata e numericamente delicata;
  \item il planning non e' un ``semplice solver balistico'', ma un selettore multi-modalita' multi-obiettivo;
  \item la separazione tra planning e fire gate e' intenzionale e corretta;
  \item i test esistenti coprono bene i contratti logici principali, ma non sostituiscono validazione hardware su TF, latenza reale e balistica reale.
\end{enumerate}

\section{Definizioni di termini tecnici, acronimi e parole ambigue}

\begin{description}[leftmargin=1.8em, style=nextline]
  \item[Armor / armatura] Piastra bersaglio del robot avversario, osservata dalla camera e usata come riferimento di tiro.
  \item[Bounding box] Rettangolo 2D che approssima la posizione dell'oggetto nell'immagine.
  \item[Callback per-frame] Funzione eseguita a ogni nuovo messaggio di detection.
  \item[Coast] Evoluzione del tracker senza misura nuova, basata solo sulla prediction del modello.
  \item[Detection] Osservazione grezza prodotta dal detector visivo.
  \item[EKF] Extended Kalman Filter, filtro di stato per sistemi non lineari linearizzati localmente.
  \item[Face jump] Cambio di faccia osservata del robot, tipicamente dovuto a una rotazione del corpo.
  \item[Face obliquity / obliquita'] Angolo tra la normale della faccia e la direzione verso la camera. Valori alti degradano la qualita' del PnP.
  \item[Fire gate] Blocco logico finale che decide se il sistema puo' davvero sparare.
  \item[Fresh] Dato ancora considerato temporalmente affidabile.
  \item[IPPE] Metodo PnP specializzato per oggetti planari.
  \item[Kalman gain] Guadagno che bilancia fiducia nel modello e fiducia nella misura.
  \item[Latency bias / time bias] Stima della latenza da compensare con una predizione in avanti.
  \item[Mahalanobis distance] Distanza statistica normalizzata dall'incertezza del filtro.
  \item[Matched face] Faccia che il tracker ritiene attualmente associata alla misura corrente.
  \item[Measurement stale] Misura troppo vecchia per essere usata in modo sicuro nel fuoco.
  \item[Observable] Faccia che, dato l'angolo rispetto alla camera, e' ancora considerata geometricamente osservabile.
  \item[Planning context] Pacchetto di parametri runtime dato al planner: posa gimbal, limiti, latenza, bias, storico target.
  \item[PnP] Perspective-n-Point, problema di stima posa da punti 2D/3D corrispondenti.
  \item[Pose source] Origine dei dati di orientazione del gimbal/camera.
  \item[Prediction horizon] Intervallo temporale nel futuro su cui si fa la previsione per il tiro.
  \item[Reachable] Soluzione balistica meccanicamente raggiungibile dal gimbal.
  \item[Scalar KF] Filtro di Kalman unidimensionale usato qui per stimare parametri quasi statici come il raggio.
  \item[Shot plan] Descrizione completa di un possibile tiro verso una specifica faccia di un target.
  \item[Slew] Movimento richiesto al gimbal per raggiungere una soluzione di tiro.
  \item[TF2] Libreria ROS per trasformazioni tra frame di riferimento.
  \item[TrackSnapshot] Fotografia immutabile dello stato di un tracker, usata dal planner.
  \item[Tracking frame] Frame stabile in cui il target viene stimato e pubblicato, di default \code{odom}.
  \item[Transport delay] Ritardo osservato tra timestamp del frame e tempo corrente del nodo.
  \item[Twist di compatibilita'] Uso non standard di \code{geometry\_msgs::msg::Twist} per trasportare yaw/pitch assoluti legacy su \code{angular.z/y}; non va interpretato come velocita' fisica.
\end{description}

\section{Fatti verificati, inferenze plausibili e ipotesi}

\subsection{Fatti verificati direttamente da codice e test}

\begin{enumerate}[leftmargin=1.4em]
  \item L'entrypoint ROS~2 e' \code{rm\_auto\_aim::AutoAimNode}, registrato come componente.
  \item Il callback principale viene attivato da \topic{/detector/armors} con tipo \code{vision\_msgs::msg::Detection2DArray}.
  \item Il solver PnP viene creato solo dopo il primo messaggio \topic{/camera_info}; prima di allora il blocco di measurement restituisce zero armature.
  \item La subscription a \topic{/camera_info} viene distrutta dopo il primo messaggio.
  \item In \code{pose\_source=micro\_pose}, il package interpreta \code{pose.position.x} come pitch e \code{pose.position.y} come yaw.
  \item In \code{pose\_source=none}, \code{GimbalPoseAdapter::isFresh()} restituisce sempre vero.
  \item Il tracker supporta concettualmente 2, 3 e 4 armature, ma l'implementazione attuale imposta sempre \code{NORMAL\_4}.
  \item L'assegnazione detection-tracker e' basata su Mahalanobis; i tracker in DETECTING sono penalizzati di un fattore \(4\).
  \item Il tracker ha quattro stati: LOST, DETECTING, TRACKING, TEMP\_LOST.
  \item Il fire gate blocca esplicitamente il fuoco quando il piano non e' in tracking valido.
  \item Il fire gate blocca esplicitamente il fuoco quando il piano e' TEMP\_LOST.
  \item Se \code{pose\_source=none}, il fuoco richiede \code{debug.allow\_fire\_without\_pose=true}.
  \item Il planner puo' produrre modalita' \code{VISIBLE\_DIRECT}, \code{PREDICTED\_DIRECT} e \code{INDIRECT}.
  \item Il \code{GimbalCmd} viene pubblicato con yaw/pitch assoluti in gradi; non in radianti.
  \item Il \code{Twist} pubblicato e' riutilizzato come compatibilita' parziale: \code{angular.y/z} portano pitch/yaw assoluti in radianti, mentre il fire flag affidabile resta \code{GimbalCmd::fire\_cmd}.
  \item I test coprono almeno: fire gate, latency compensator, engagement planner, gimbal pose adapter, gimbal command controller, tracker runtime factory e tracker runtime manager.
  \item Nel nodo attuale i pesi della cost function vengono dichiarati come parametri ROS \code{cost.*} e possono essere sovrascritti dal launch.
  \item Il PnP ha un gate esplicito di errore di reproiezione tramite \code{pnp.max\_reprojection\_error}.
\end{enumerate}

\subsection{Inferenze plausibili, forti ma non totalmente dimostrate dal solo codice}

\begin{enumerate}[leftmargin=1.4em]
  \item Il package e' pensato come consolidamento operativo dei package legacy \code{armor\_tracker} e \code{rm\_trajectory}.
  \item Il topic \topic{/cmd_vel} viene mantenuto soprattutto per retrocompatibilita' con consumatori legacy a valle.
  \item La semantica pratica di \code{target\_classes} nel progetto sembra essere piu' vicina a ``filtrare colore/team'' che a distinguere il tipo robotico reale.
  \item La modalita' \code{pose\_source=none} e' pensata primariamente per debug, bag replay o ambienti dove la camera e' trattata come orientata in modo noto.
  \item Il README del package e l'owner-path launch indicano che il flusso voluto di sviluppo si sta spostando sul package unificato, mentre altri launch del repository possono ancora riflettere l'architettura separata precedente.
\end{enumerate}

\subsection{Ipotesi, limiti epistemici e punti che il codice da solo non puo' provare}

\begin{enumerate}[leftmargin=1.4em]
  \item Il codice non puo' dimostrare da solo la precisione balistica reale sul robot fisico: servono log, test di tiro e calibrazione.
  \item Il codice non puo' garantire che i timestamp dei frame e dell'IMU siano coerenti a livello hardware; questo incide fortemente su \code{time\_bias} e indirect mode.
  \item L'effetto pratico della soglia \code{pnp.max\_reprojection\_error} dipende dalla qualita' del detector, dalla calibrazione camera e dalla precisione della bbox.
  \item Non e' possibile dedurre dal solo package se i segni \code{gimbal.yaw\_sign} e \code{gimbal.pitch\_sign} siano corretti per ogni montaggio meccanico.
  \item Le prestazioni su target a 2 o 3 facce non sono verificabili nel runtime attuale, perche' il tracker non sta inferendo davvero quei casi.
\end{enumerate}

\section{Indicazioni pratiche: cosa fare, come verificare, errori da evitare}

\subsection{Cosa fare per integrare o capire il package nel modo corretto}

\begin{enumerate}[leftmargin=1.4em]
  \item Verificare prima di tutto \topic{/camera_info}. Senza intrinseci non esistono misure 3D vere.
  \item Verificare il TF \code{target\_frame -> camera\_color\_optical\_frame}. Se quel TF e' sbagliato, tutto il tracking e il planning partono gia' storti.
  \item Verificare che \topic{/detector/armors} abbia classi coerenti con \code{target\_classes}.
  \item Verificare il comportamento di \topic{/tracker/info} prima di toccare la balistica. Se measurement e tracking sono sporchi, la balistica non puo' salvarli.
  \item Solo dopo measurement e tracking iniziare a toccare \code{time\_bias}, \code{angular\_window}, \code{indirect\_*}, \code{cost.*} e i clamp dei comandi.
\end{enumerate}

\subsection{Come verificare il sistema in pratica}

\paragraph{Prima verifica: build e indice C++.}
\begin{lstlisting}[style=shellstyle,language=bash,caption={Build minimo e compile database per editor}]
source /opt/ros/humble/setup.bash
colcon build --packages-up-to auto_aim_targeting --cmake-args -DCMAKE_EXPORT_COMPILE_COMMANDS=ON
ln -sf build/compile_commands.json compile_commands.json
\end{lstlisting}

Se la build passa ma l'IDE mostra include mancanti, il problema e' quasi sempre l'indice C++ dell'editor, non il codice. Il file \code{compile\_commands.json} contiene gli include reali usati dal compilatore: ROS Humble, Eigen, OpenCV, \pkg\ e \code{auto\_aim\_interfaces}.

\paragraph{Verifica runtime.}
\begin{lstlisting}[style=shellstyle,language=bash,caption={Comandi minimi utili per ispezionare il runtime}]
ros2 topic echo /tracker/target
ros2 topic echo /tracker/targets
ros2 topic echo /tracker/info
ros2 topic echo /tracker/cmd_gimbal
ros2 topic echo /detections_output/optimal_target
ros2 service call /tracker/reset std_srvs/srv/Trigger "{}"
\end{lstlisting}

Checklist interpretativa:
\begin{enumerate}[leftmargin=1.4em]
  \item Se \topic{/tracker/target} resta sempre vuoto ma il detector vede oggetti, controllare \topic{/camera_info} e TF.
  \item Se il target compare ma \code{tracking=false} a lungo, controllare \code{tracker.tracking\_thres}, associazione e qualità misura.
  \item Se il target segue bene ma non spara, leggere il \code{blocker} del fire gate nel codice o strumentarlo esplicitamente nei log.
  \item Se i marker delle armature sono geometricamente sfasati, partire da \code{light\_ratio}, offset camera e convenzioni di pose source.
  \item Se il sistema spara sistematicamente dietro o davanti su target in moto, toccare \code{time\_bias} solo dopo aver validato timestamp e \code{bullet\_speed}.
\end{enumerate}

\subsection{Ordine corretto di debug}

\begin{enumerate}[leftmargin=1.4em]
  \item \textbf{Detector}: le bbox sono giuste? Le classi sono giuste?
  \item \textbf{Measurement}: la posa 3D e' ragionevole? La distanza e lo yaw hanno senso?
  \item \textbf{TF / pose source}: il target e' nel frame giusto? Gli assi non sono invertiti?
  \item \textbf{Tracking}: il target resta coerente durante rotazione e occlusioni brevi?
  \item \textbf{Planning}: cambia correttamente tra visible, predicted e indirect?
  \item \textbf{Fire gate}: il blocco del fuoco e' quello atteso?
  \item \textbf{Command output}: clamp, unita' di misura e convenzione assoluta degli angoli sono compatibili con l'attuatore reale?
\end{enumerate}

\subsection{Errori concettuali da evitare}

\begin{enumerate}[leftmargin=1.4em]
  \item Non trattare \code{light\_ratio} come un semplice filtro estetico: cambia direttamente la geometria data al PnP.
  \item Non assumere che \code{pose\_source=none} sia adatto al robot reale solo perche' il package continua a funzionare.
  \item Non interpretare \topic{/cmd_vel} come vera twist ROS standard.
  \item Non credere che il runtime attuale distingua davvero robot a 2, 3 e 4 armature: oggi assume sempre il caso a 4.
  \item Non usare \code{cost.*} per correggere errori a monte: i pesi sono esposti, ma non possono compensare TF, timestamp o PnP sbagliati.
  \item Non correggere errori di measurement con parametri di planning: e' il modo piu' rapido per ottenere un sistema fragile.
  \item Non ignorare la differenza tra \code{valid} e \code{reachable}: un piano puo' essere numericamente valido ma meccanicamente impossibile.
\end{enumerate}

\subsection{Messaggi finali davvero utili da ricordare}

\begin{enumerate}[leftmargin=1.4em]
  \item \pkg\ e' un package di \emph{state estimation + decision making}, non un semplice trasformatore di detection in comando.
  \item La qualità del tiro dipende piu' dalla qualita' di measurement, timestamps e TF che da piccoli ritocchi ai pesi del planner.
  \item Il concetto chiave del tracker e' ``modello del robot'', non ``inseguimento della faccia''.
  \item Il concetto chiave del planner e' ``sincronizzare geometria, latenza e balistica'', non semplicemente ``puntare verso il target''.
\end{enumerate}

\end{document}
